\documentclass[12pt, ukrainian]{article}
\usepackage[a4paper]{geometry}
\setlength\parindent{0mm}
\geometry{top=1.1cm,bottom=1.1cm,left=1.1cm,right=1.1cm}
\pagestyle{empty}
\usepackage[T2A]{fontenc}
\usepackage[utf8]{inputenc}
\usepackage[ukrainian]{babel}
\usepackage{graphicx}
\usepackage{caption}
\DeclareCaptionFormat{nocaption}{}
\captionsetup{format=nocaption,aboveskip=0pt,belowskip=0pt}
\usepackage[Export]{adjustbox}
\adjustboxset{max size={0.9\linewidth}{0.9\paperheight}}
\begin{document}
{{22.11.2024 Контрольна робота \No 2, варіант  \No \textbf { 1 }\par}}
{
%\begin {large}

Тестова частина. Завдання 1-5: Що буде виведено за виконання (повідомлення інтерпретатора про помилку позначати error)?

Задачі (на звороті): написати функцію для розв'язання задачі.

\vspace{5mm} \begin{minipage}[t]{8.5cm}
               
{\begin{center}Завдання {1}\end{center}}
\vspace{-6mm}

\begin{verbatim}
a,b,c=4,5,9
def g(b):
    global c
    a=2
    b*=4
    c=3
    return a+b+c

a,b,c=1,8,0
print(g(b),a,b,c)
\end{verbatim}


\vspace{5mm}
               
{\begin{center}Завдання {2}\end{center}}
\vspace{-6mm}

\begin{verbatim}
def f():
    try:
        res = int(8//2)
        return 43
    except TypeError: return 9
    except Exception: return 3
    return res

print(f())
\end{verbatim}


\vspace{5mm}
               
{\begin{center}Завдання {3}\end{center}}
\vspace{-6mm}

\begin{verbatim}
j=8
while j<=10:
    j+=1
print(j)
\end{verbatim}


\end{minipage}
\vline \hspace{6mm}
\begin{minipage}[t]{10.7cm}
               
{\begin{center}Завдання {4}\end{center}}
\vspace{-6mm}

\begin{verbatim}
for a in range(2,2+4,-2):
    if a>=6:
        break
    print(a, end=' ')
    a=4
else:
    print(13, end=' ')
print(a, end=' ')
\end{verbatim}


\vspace{5mm}
               
{\begin{center}Завдання {5}\end{center}}
\vspace{-6mm}

\begin{verbatim}
def f(n):
    if n > 9:
        return f(n // 100) + 1
    else:
        return 1

print(f(123456))
\end{verbatim}


\end{minipage}

\newpage

{{22.11.2024 Контрольна робота \No 2, варіант  \No \textbf { 2 }\par}}
{
%\begin {large}

Тестова частина. Завдання 1-5: Що буде виведено за виконання (повідомлення інтерпретатора про помилку позначати error)?

Задачі (на звороті): написати функцію для розв'язання задачі.

\vspace{5mm} \begin{minipage}[t]{8.5cm}
               
{\begin{center}Завдання {1}\end{center}}
\vspace{-6mm}

\begin{verbatim}
a,b,c=9,0,5
def h(b):
    a=3
    b=4
    c=1
    return a+b+c

a,b,c=1,3,7
print(h(a),a,b,c)
\end{verbatim}


\vspace{5mm}
               
{\begin{center}Завдання {2}\end{center}}
\vspace{-6mm}

\begin{verbatim}
def f():
    try:
        res = 7>=7
    except KeyboardInterrupt: return 4
    except ValueError: return 0
    else: return 33
    finally: return 22
    return res

print(f())
\end{verbatim}


\vspace{5mm}
               
{\begin{center}Завдання {3}\end{center}}
\vspace{-6mm}

\begin{verbatim}
i=3
while i<12:
    i+=2
print(i)
\end{verbatim}


\end{minipage}
\vline \hspace{6mm}
\begin{minipage}[t]{10.7cm}
               
{\begin{center}Завдання {4}\end{center}}
\vspace{-6mm}

\begin{verbatim}
for c in range(-10,-10+5,-3):
    if c<=5:
        pass
    print(c, end=' ')
    c=6
    if c<=6:
        break
else:
    print(c, end=' ')
print(c, end=' ')
\end{verbatim}


\vspace{5mm}
               
{\begin{center}Завдання {5}\end{center}}
\vspace{-6mm}

\begin{verbatim}
def f(n):
    print(n % 10, end=' ')
    if n > 9:
        f(n // 10)
    

f(543216)
\end{verbatim}


\end{minipage}

\newpage

{{22.11.2024 Контрольна робота \No 2, варіант  \No \textbf { 3 }\par}}
{
%\begin {large}

Тестова частина. Завдання 1-5: Що буде виведено за виконання (повідомлення інтерпретатора про помилку позначати error)?

Задачі (на звороті): написати функцію для розв'язання задачі.

\vspace{5mm} \begin{minipage}[t]{8.5cm}
               
{\begin{center}Завдання {1}\end{center}}
\vspace{-6mm}

\begin{verbatim}
a,b,c=3,7,6
def f(a):
    a=1
    b-=5
    c=3
    return a+b+c

a,b,c=2,9,3
print(f(b),a,b,c)
\end{verbatim}


\vspace{5mm}
               
{\begin{center}Завдання {2}\end{center}}
\vspace{-6mm}

\begin{verbatim}
def f():
    try:
        res = int("d1")
        return 41
    except BaseException: return 2
    except ZeroDivisionError: return 5
    finally: return 24
    return res

print(f())
\end{verbatim}


\vspace{5mm}
               
{\begin{center}Завдання {3}\end{center}}
\vspace{-6mm}

\begin{verbatim}
k=15
while k>5:
    k+=-2
print(k)
\end{verbatim}


\end{minipage}
\vline \hspace{6mm}
\begin{minipage}[t]{10.7cm}
               
{\begin{center}Завдання {4}\end{center}}
\vspace{-6mm}

\begin{verbatim}
for d in range(-7,-7+3,2):
    if d>8:
        continue
    print(d, end=' ')
    d=-5
else:
    print(d, end=' ')
print(d, end=' ')
\end{verbatim}


\vspace{5mm}
               
{\begin{center}Завдання {5}\end{center}}
\vspace{-6mm}

\begin{verbatim}
def f(n):
    if n > 9:
        f(n // 100)
    print(n % 10, end=' ')
    

f(123456)
\end{verbatim}


\end{minipage}

\newpage

{{22.11.2024 Контрольна робота \No 2, варіант  \No \textbf { 4 }\par}}
{
%\begin {large}

Тестова частина. Завдання 1-5: Що буде виведено за виконання (повідомлення інтерпретатора про помилку позначати error)?

Задачі (на звороті): написати функцію для розв'язання задачі.

\vspace{5mm} \begin{minipage}[t]{8.5cm}
               
{\begin{center}Завдання {1}\end{center}}
\vspace{-6mm}

\begin{verbatim}
a,b,c=5,2,6
def h(b):
    a+=2
    b=1
    c=4
    return a+b+c

a,b,c=0,7,1
print(h(a),a,b,c)
\end{verbatim}


\vspace{5mm}
               
{\begin{center}Завдання {2}\end{center}}
\vspace{-6mm}

\begin{verbatim}
def f():
    try:
        res = int(6/0.0)
    except KeyboardInterrupt: return 2
    except TypeError: return 0
    else: return 30
    return res

print(f())
\end{verbatim}


\vspace{5mm}
               
{\begin{center}Завдання {3}\end{center}}
\vspace{-6mm}

\begin{verbatim}
n=11
while n>=4:
    n+=-3
print(n)
\end{verbatim}


\end{minipage}
\vline \hspace{6mm}
\begin{minipage}[t]{10.7cm}
               
{\begin{center}Завдання {4}\end{center}}
\vspace{-6mm}

\begin{verbatim}
for f in range(8,8+6,3):
    if f<4:
        continue
    print(f, end=' ')
    f=5
else:
    print(f, end=' ')
print(f, end=' ')
\end{verbatim}


\vspace{5mm}
               
{\begin{center}Завдання {5}\end{center}}
\vspace{-6mm}

\begin{verbatim}
def f(n):
    if n > 9:
        f(n // 10)
    else:
        print(n % 10, end=' ')

f(543216)
\end{verbatim}


\end{minipage}

\newpage

{{22.11.2024 Контрольна робота \No 2, варіант  \No \textbf { 5 }\par}}
{
%\begin {large}

Тестова частина. Завдання 1-5: Що буде виведено за виконання (повідомлення інтерпретатора про помилку позначати error)?

Задачі (на звороті): написати функцію для розв'язання задачі.

\vspace{5mm} \begin{minipage}[t]{8.5cm}
               
{\begin{center}Завдання {1}\end{center}}
\vspace{-6mm}

\begin{verbatim}
a,b,c=4,8,9
def g(a):
    global c
    a=5
    b+=3
    c=5
    return a+b+c

a,b,c=6,2,3
print(g(a),a,b,c)
\end{verbatim}


\vspace{5mm}
               
{\begin{center}Завдання {2}\end{center}}
\vspace{-6mm}

\begin{verbatim}
def f():
    try:
        res = int("9")
    except ValueError: return 1
    except ZeroDivisionError: return 4
    else: return 32
    return res

print(f())
\end{verbatim}


\vspace{5mm}
               
{\begin{center}Завдання {3}\end{center}}
\vspace{-6mm}

\begin{verbatim}
m=7
while m>=3:
    m+=-3
print(m)
\end{verbatim}


\end{minipage}
\vline \hspace{6mm}
\begin{minipage}[t]{10.7cm}
               
{\begin{center}Завдання {4}\end{center}}
\vspace{-6mm}

\begin{verbatim}
for e in range(5,5+2,-3):
    if e>7:
        break
    print(e, end=' ')
    e=-6
else:
    print(e, end=' ')
print(e, end=' ')
\end{verbatim}


\vspace{5mm}
               
{\begin{center}Завдання {5}\end{center}}
\vspace{-6mm}

\begin{verbatim}
def f(n):
    if n > 9:
            f(n // 100)
    print(n % 10, end=' ')
    

f(987654)
\end{verbatim}


\end{minipage}

\newpage

{{22.11.2024 Контрольна робота \No 2, варіант  \No \textbf { 6 }\par}}
{
%\begin {large}

Тестова частина. Завдання 1-5: Що буде виведено за виконання (повідомлення інтерпретатора про помилку позначати error)?

Задачі (на звороті): написати функцію для розв'язання задачі.

\vspace{5mm} \begin{minipage}[t]{8.5cm}
               
{\begin{center}Завдання {1}\end{center}}
\vspace{-6mm}

\begin{verbatim}
a,b,c=4,2,8
def f(a):
    global c
    a=5
    b=2
    c=5
    return a+b+c

a,b,c=6,0,6
print(f(b),a,b,c)
\end{verbatim}


\vspace{5mm}
               
{\begin{center}Завдання {2}\end{center}}
\vspace{-6mm}

\begin{verbatim}
def f():
    try:
        res = int(6%0)
        return 44
    except BaseException: return 7
    except Exception: return 3
    finally: return 23
    return res

print(f())
\end{verbatim}


\vspace{5mm}
               
{\begin{center}Завдання {3}\end{center}}
\vspace{-6mm}

\begin{verbatim}
m=7
while m<7:
    m+=2
print(m)
\end{verbatim}


\end{minipage}
\vline \hspace{6mm}
\begin{minipage}[t]{10.7cm}
               
{\begin{center}Завдання {4}\end{center}}
\vspace{-6mm}

\begin{verbatim}
for b in range(3,3+5,-2):
    if b<3:
        continue
    print(b, end=' ')
    b=9
else:
    print(b, end=' ')
print(b, end=' ')
\end{verbatim}


\vspace{5mm}
               
{\begin{center}Завдання {5}\end{center}}
\vspace{-6mm}

\begin{verbatim}
def f(n):
    print(n % 10, end=' ')
    if n > 9:
        f(n // 10)
    

f(123456)
\end{verbatim}


\end{minipage}

\newpage

{{22.11.2024 Контрольна робота \No 2, варіант  \No \textbf { 7 }\par}}
{
%\begin {large}

Тестова частина. Завдання 1-5: Що буде виведено за виконання (повідомлення інтерпретатора про помилку позначати error)?

Задачі (на звороті): написати функцію для розв'язання задачі.

\vspace{5mm} \begin{minipage}[t]{8.5cm}
               
{\begin{center}Завдання {1}\end{center}}
\vspace{-6mm}

\begin{verbatim}
a,b,c=0,5,1
def h(b):
    a=4
    b-=2
    c=1
    return a+b+c

a,b,c=4,7,8
print(h(b),a,b,c)
\end{verbatim}


\vspace{5mm}
               
{\begin{center}Завдання {2}\end{center}}
\vspace{-6mm}

\begin{verbatim}
def f():
    try:
        res = int("5")
    except Exception: return 8
    except BaseException: return 5
    finally: return 21
    return res

print(f())
\end{verbatim}


\vspace{5mm}
               
{\begin{center}Завдання {3}\end{center}}
\vspace{-6mm}

\begin{verbatim}
t=12
while t>2:
    t+=-2
print(t)
\end{verbatim}


\end{minipage}
\vline \hspace{6mm}
\begin{minipage}[t]{10.7cm}
               
{\begin{center}Завдання {4}\end{center}}
\vspace{-6mm}

\begin{verbatim}
for c in range(9,9+6,3):
    if c<=4:
        pass
    print(c, end=' ')
    c=2
else:
    print(c, end=' ')
print(c, end=' ')
\end{verbatim}


\vspace{5mm}
               
{\begin{center}Завдання {5}\end{center}}
\vspace{-6mm}

\begin{verbatim}
def f(n):
    if n > 9:
        f(n // 10)
    print(n % 10, end=' ')
    

f(123456)
\end{verbatim}


\end{minipage}

\newpage

{{22.11.2024 Контрольна робота \No 2, варіант  \No \textbf { 8 }\par}}
{
%\begin {large}

Тестова частина. Завдання 1-5: Що буде виведено за виконання (повідомлення інтерпретатора про помилку позначати error)?

Задачі (на звороті): написати функцію для розв'язання задачі.

\vspace{5mm} \begin{minipage}[t]{8.5cm}
               
{\begin{center}Завдання {1}\end{center}}
\vspace{-6mm}

\begin{verbatim}
a,b,c=4,9,1
def f(a):
    global c
    a*=2
    b=4
    c=5
    return a+b+c

a,b,c=0,5,2
print(f(a),a,b,c)
\end{verbatim}


\vspace{5mm}
               
{\begin{center}Завдання {2}\end{center}}
\vspace{-6mm}

\begin{verbatim}
def f():
    try:
        res = 8<0
        return 42
    except ValueError: return 6
    except KeyboardInterrupt: return 3
    else: return 31
    return res

print(f())
\end{verbatim}


\vspace{5mm}
               
{\begin{center}Завдання {3}\end{center}}
\vspace{-6mm}

\begin{verbatim}
j=8
while j>0:
    j+=-2
print(j)
\end{verbatim}


\end{minipage}
\vline \hspace{6mm}
\begin{minipage}[t]{10.7cm}
               
{\begin{center}Завдання {4}\end{center}}
\vspace{-6mm}

\begin{verbatim}
for b in range(-1,-1+2,2):
    if b>=6:
        continue
    print(b, end=' ')
    b=-2
else:
    print(b, end=' ')
print(b, end=' ')
\end{verbatim}


\vspace{5mm}
               
{\begin{center}Завдання {5}\end{center}}
\vspace{-6mm}

\begin{verbatim}
def f(n):
    if n > 2:
        f(n // 2)
    else:
        print(n % 2, end=' ')

f(16)
\end{verbatim}


\end{minipage}

\newpage

{{22.11.2024 Контрольна робота \No 2, варіант  \No \textbf { 9 }\par}}
{
%\begin {large}

Тестова частина. Завдання 1-5: Що буде виведено за виконання (повідомлення інтерпретатора про помилку позначати error)?

Задачі (на звороті): написати функцію для розв'язання задачі.

\vspace{5mm} \begin{minipage}[t]{8.5cm}
               
{\begin{center}Завдання {1}\end{center}}
\vspace{-6mm}

\begin{verbatim}
a,b,c=7,6,3
def f(b):
    a-=3
    b=1
    c=2
    return a+b+c

a,b,c=8,6,2
print(f(a),a,b,c)
\end{verbatim}


\vspace{5mm}
               
{\begin{center}Завдання {2}\end{center}}
\vspace{-6mm}

\begin{verbatim}
def f():
    try:
        res = int(1%1)
        return 40
    except ZeroDivisionError: return 2
    except TypeError: return 9
    return res

print(f())
\end{verbatim}


\vspace{5mm}
               
{\begin{center}Завдання {3}\end{center}}
\vspace{-6mm}

\begin{verbatim}
i=13
while i>=9:
    i+=-3
print(i)
\end{verbatim}


\end{minipage}
\vline \hspace{6mm}
\begin{minipage}[t]{10.7cm}
               
{\begin{center}Завдання {4}\end{center}}
\vspace{-6mm}

\begin{verbatim}
for f in range(-2,-2+4,-3):
    if f>=8:
        continue
    print(f, end=' ')
    f=-1
else:
    print(f, end=' ')
print(f, end=' ')
\end{verbatim}


\vspace{5mm}
               
{\begin{center}Завдання {5}\end{center}}
\vspace{-6mm}

\begin{verbatim}
def f(n):
    if n > 9:
        return f(n // 10) + 1
    else:
        return 1

print(f(123456))
\end{verbatim}


\end{minipage}

\newpage

{{22.11.2024 Контрольна робота \No 2, варіант  \No \textbf { 10 }\par}}
{
%\begin {large}

Тестова частина. Завдання 1-5: Що буде виведено за виконання (повідомлення інтерпретатора про помилку позначати error)?

Задачі (на звороті): написати функцію для розв'язання задачі.

\vspace{5mm} \begin{minipage}[t]{8.5cm}
               
{\begin{center}Завдання {1}\end{center}}
\vspace{-6mm}

\begin{verbatim}
a,b,c=7,9,0
def g(a):
    global c
    a=4
    b+=1
    c=3
    return a+b+c

a,b,c=1,5,8
print(g(a),a,b,c)
\end{verbatim}


\vspace{5mm}
               
{\begin{center}Завдання {2}\end{center}}
\vspace{-6mm}

\begin{verbatim}
def f():
    try:
        res = int("a0")
    except BaseException: return 4
    except TypeError: return 7
    else: return 34
    finally: return 20
    return res

print(f())
\end{verbatim}


\vspace{5mm}
               
{\begin{center}Завдання {3}\end{center}}
\vspace{-6mm}

\begin{verbatim}
k=4
while k<=5:
    k+=1
print(k)
\end{verbatim}


\end{minipage}
\vline \hspace{6mm}
\begin{minipage}[t]{10.7cm}
               
{\begin{center}Завдання {4}\end{center}}
\vspace{-6mm}

\begin{verbatim}
for a in range(6,6+3,3):
    if a<7:
        pass
    print(a, end=' ')
    a=-3
else:
    print(13, end=' ')
print(a, end=' ')
\end{verbatim}


\vspace{5mm}
               
{\begin{center}Завдання {5}\end{center}}
\vspace{-6mm}

\begin{verbatim}
def f(n):
    print(n % 2, end=' ')
    if n > 2:
        f(n // 2)
    

f(16)
\end{verbatim}


\end{minipage}

\newpage

{{22.11.2024 Контрольна робота \No 2, варіант  \No \textbf { 11 }\par}}
{
%\begin {large}

Тестова частина. Завдання 1-5: Що буде виведено за виконання (повідомлення інтерпретатора про помилку позначати error)?

Задачі (на звороті): написати функцію для розв'язання задачі.

\vspace{5mm} \begin{minipage}[t]{8.5cm}
               
{\begin{center}Завдання {1}\end{center}}
\vspace{-6mm}

\begin{verbatim}
a,b,c=4,3,0
def h(a):
    a=5
    b*=5
    c=2
    return a+b+c

a,b,c=1,8,3
print(h(b),a,b,c)
\end{verbatim}


\vspace{5mm}
               
{\begin{center}Завдання {2}\end{center}}
\vspace{-6mm}

\begin{verbatim}
def f():
    try:
        res = int("9")
        return 45
    except KeyboardInterrupt: return 8
    except ValueError: return 2
    return res

print(f())
\end{verbatim}


\vspace{5mm}
               
{\begin{center}Завдання {3}\end{center}}
\vspace{-6mm}

\begin{verbatim}
t=2
while t<6:
    t+=1
print(t)
\end{verbatim}


\end{minipage}
\vline \hspace{6mm}
\begin{minipage}[t]{10.7cm}
               
{\begin{center}Завдання {4}\end{center}}
\vspace{-6mm}

\begin{verbatim}
for d in range(-4,-4+5,-2):
    if d<=5:
        break
    print(d, end=' ')
    d=0
else:
    print(d, end=' ')
print(d, end=' ')
\end{verbatim}


\vspace{5mm}
               
{\begin{center}Завдання {5}\end{center}}
\vspace{-6mm}

\begin{verbatim}
def f(n):
    if n > 2:
        f(n // 2)
    print(n % 2, end=' ')
    

f(16)
\end{verbatim}


\end{minipage}

\newpage

{{22.11.2024 Контрольна робота \No 2, варіант  \No \textbf { 12 }\par}}
{
%\begin {large}

Тестова частина. Завдання 1-5: Що буде виведено за виконання (повідомлення інтерпретатора про помилку позначати error)?

Задачі (на звороті): написати функцію для розв'язання задачі.

\vspace{5mm} \begin{minipage}[t]{8.5cm}
               
{\begin{center}Завдання {1}\end{center}}
\vspace{-6mm}

\begin{verbatim}
a,b,c=3,8,1
def g(a):
    a=1
    b=2
    c=3
    return a+b+c

a,b,c=5,7,4
print(g(a),a,b,c)
\end{verbatim}


\vspace{5mm}
               
{\begin{center}Завдання {2}\end{center}}
\vspace{-6mm}

\begin{verbatim}
def f():
    try:
        res = int(0/0.0)
    except Exception: return 5
    except ZeroDivisionError: return 4
    else: return 35
    finally: return 25
    return res

print(f())
\end{verbatim}


\vspace{5mm}
               
{\begin{center}Завдання {3}\end{center}}
\vspace{-6mm}

\begin{verbatim}
m=14
while m>=6:
    m+=-2
print(m)
\end{verbatim}


\end{minipage}
\vline \hspace{6mm}
\begin{minipage}[t]{10.7cm}
               
{\begin{center}Завдання {4}\end{center}}
\vspace{-6mm}

\begin{verbatim}
for e in range(0,0+3,2):
    if e>3:
        continue
    print(e, end=' ')
    e=-7
else:
    print(e, end=' ')
print(e, end=' ')
\end{verbatim}


\vspace{5mm}
               
{\begin{center}Завдання {5}\end{center}}
\vspace{-6mm}

\begin{verbatim}
def f(n):
    if n <= 9:
        print(n % 10, end=' ')
    else:
        f(n // 10)

f(987651)
\end{verbatim}


\end{minipage}

\newpage

{{22.11.2024 Контрольна робота \No 2, варіант  \No \textbf { 13 }\par}}
{
%\begin {large}

Тестова частина. Завдання 1-5: Що буде виведено за виконання (повідомлення інтерпретатора про помилку позначати error)?

Задачі (на звороті): написати функцію для розв'язання задачі.

\vspace{5mm} \begin{minipage}[t]{8.5cm}
               
{\begin{center}Завдання {1}\end{center}}
\vspace{-6mm}

\begin{verbatim}
a,b,c=7,2,4
def h(b):
    global c
    a*=1
    b=4
    c=3
    return a+b+c

a,b,c=5,6,9
print(h(a),a,b,c)
\end{verbatim}


\vspace{5mm}
               
{\begin{center}Завдання {2}\end{center}}
\vspace{-6mm}

\begin{verbatim}
def f():
    try:
        res = int(7//2)
    except KeyboardInterrupt: return 3
    except Exception: return 6
    return res

print(f())
\end{verbatim}


\vspace{5mm}
               
{\begin{center}Завдання {3}\end{center}}
\vspace{-6mm}

\begin{verbatim}
n=1
while n<=13:
    n+=2
print(n)
\end{verbatim}


\end{minipage}
\vline \hspace{6mm}
\begin{minipage}[t]{10.7cm}
               
{\begin{center}Завдання {4}\end{center}}
\vspace{-6mm}

\begin{verbatim}
for b in range(-8,-8+6,-3):
    if b>=7:
        continue
    print(b, end=' ')
    b=-10
else:
    print(b, end=' ')
print(b, end=' ')
\end{verbatim}


\vspace{5mm}
               
{\begin{center}Завдання {5}\end{center}}
\vspace{-6mm}

\begin{verbatim}
def f(n):
    if n > 9:
        f(n // 100)
    else:
        print(n % 10, end=' ')

f(987654)
\end{verbatim}


\end{minipage}

\newpage

{{22.11.2024 Контрольна робота \No 2, варіант  \No \textbf { 14 }\par}}
{
%\begin {large}

Тестова частина. Завдання 1-5: Що буде виведено за виконання (повідомлення інтерпретатора про помилку позначати error)?

Задачі (на звороті): написати функцію для розв'язання задачі.

\vspace{5mm} \begin{minipage}[t]{8.5cm}
               
{\begin{center}Завдання {1}\end{center}}
\vspace{-6mm}

\begin{verbatim}
a,b,c=2,7,6
def g(b):
    global c
    a-=4
    b=3
    c=2
    return a+b+c

a,b,c=1,5,8
print(g(b),a,b,c)
\end{verbatim}


\vspace{5mm}
               
{\begin{center}Завдання {2}\end{center}}
\vspace{-6mm}

\begin{verbatim}
def f():
    try:
        res = int("b1")
        return 42
    except BaseException: return 5
    except ValueError: return 4
    else: return 33
    finally: return 24
    return res

print(f())
\end{verbatim}


\vspace{5mm}
               
{\begin{center}Завдання {3}\end{center}}
\vspace{-6mm}

\begin{verbatim}
i=9
while i>1:
    i+=-3
print(i)
\end{verbatim}


\end{minipage}
\vline \hspace{6mm}
\begin{minipage}[t]{10.7cm}
               
{\begin{center}Завдання {4}\end{center}}
\vspace{-6mm}

\begin{verbatim}
for e in range(1,1+4,-2):
    if e>8:
        continue
    print(e, end=' ')
    e=1
    if e>8:
        break
else:
    print(e, end=' ')
print(e, end=' ')
\end{verbatim}


\vspace{5mm}
               
{\begin{center}Завдання {5}\end{center}}
\vspace{-6mm}

\begin{verbatim}
def f(n):
    if n > 9:
        f(n // 10)
    print(n % 10, end=' ')
    

f(543216)
\end{verbatim}


\end{minipage}

\newpage

{{22.11.2024 Контрольна робота \No 2, варіант  \No \textbf { 15 }\par}}
{
%\begin {large}

Тестова частина. Завдання 1-5: Що буде виведено за виконання (повідомлення інтерпретатора про помилку позначати error)?

Задачі (на звороті): написати функцію для розв'язання задачі.

\vspace{5mm} \begin{minipage}[t]{8.5cm}
               
{\begin{center}Завдання {1}\end{center}}
\vspace{-6mm}

\begin{verbatim}
a,b,c=3,9,0
def f(a):
    a+=1
    b=5
    c=4
    return a+b+c

a,b,c=4,4,2
print(f(b),a,b,c)
\end{verbatim}


\vspace{5mm}
               
{\begin{center}Завдання {2}\end{center}}
\vspace{-6mm}

\begin{verbatim}
def f():
    try:
        res = 7==3
    except TypeError: return 0
    except ZeroDivisionError: return 2
    else: return 34
    finally: return 21
    return res

print(f())
\end{verbatim}


\vspace{5mm}
               
{\begin{center}Завдання {3}\end{center}}
\vspace{-6mm}

\begin{verbatim}
j=6
while j>7:
    j+=-2
print(j)
\end{verbatim}


\end{minipage}
\vline \hspace{6mm}
\begin{minipage}[t]{10.7cm}
               
{\begin{center}Завдання {4}\end{center}}
\vspace{-6mm}

\begin{verbatim}
for d in range(10,10+2,3):
    if d<=5:
        pass
    print(d, end=' ')
    d=-4
else:
    print(d, end=' ')
print(d, end=' ')
\end{verbatim}


\vspace{5mm}
               
{\begin{center}Завдання {5}\end{center}}
\vspace{-6mm}

\begin{verbatim}
def f(n):
    if n <= 9:
        print(n % 10, end=' ')
    else:
        f(n // 10)

f(543216)
\end{verbatim}


\end{minipage}

\newpage

{{22.11.2024 Контрольна робота \No 2, варіант  \No \textbf { 16 }\par}}
{
%\begin {large}

Тестова частина. Завдання 1-5: Що буде виведено за виконання (повідомлення інтерпретатора про помилку позначати error)?

Задачі (на звороті): написати функцію для розв'язання задачі.

\vspace{5mm} \begin{minipage}[t]{8.5cm}
               
{\begin{center}Завдання {1}\end{center}}
\vspace{-6mm}

\begin{verbatim}
a,b,c=2,9,9
def f(b):
    global c
    a=4
    b=2
    c=4
    return a+b+c

a,b,c=0,5,8
print(f(b),a,b,c)
\end{verbatim}


\vspace{5mm}
               
{\begin{center}Завдання {2}\end{center}}
\vspace{-6mm}

\begin{verbatim}
def f():
    try:
        res = 3<=5
        return 44
    except KeyboardInterrupt: return 8
    except ValueError: return 9
    return res

print(f())
\end{verbatim}


\vspace{5mm}
               
{\begin{center}Завдання {3}\end{center}}
\vspace{-6mm}

\begin{verbatim}
k=0
while k<8:
    k+=1
print(k)
\end{verbatim}


\end{minipage}
\vline \hspace{6mm}
\begin{minipage}[t]{10.7cm}
               
{\begin{center}Завдання {4}\end{center}}
\vspace{-6mm}

\begin{verbatim}
for c in range(7,7+2,2):
    if c<3:
        break
    print(c, end=' ')
    c=-9
else:
    print(c, end=' ')
print(c, end=' ')
\end{verbatim}


\vspace{5mm}
               
{\begin{center}Завдання {5}\end{center}}
\vspace{-6mm}

\begin{verbatim}
def f(n):
    print(n % 10, end=' ')
    if n > 9:
           f(n // 100)
    

f(123456)
\end{verbatim}


\end{minipage}

\newpage

{{22.11.2024 Контрольна робота \No 2, варіант  \No \textbf { 17 }\par}}
{
%\begin {large}

Тестова частина. Завдання 1-5: Що буде виведено за виконання (повідомлення інтерпретатора про помилку позначати error)?

Задачі (на звороті): написати функцію для розв'язання задачі.

\vspace{5mm} \begin{minipage}[t]{8.5cm}
               
{\begin{center}Завдання {1}\end{center}}
\vspace{-6mm}

\begin{verbatim}
a,b,c=1,2,3
def g(a):
    global c
    a=5
    b=3
    c=1
    return a+b+c

a,b,c=6,7,4
print(g(b),a,b,c)
\end{verbatim}


\vspace{5mm}
               
{\begin{center}Завдання {2}\end{center}}
\vspace{-6mm}

\begin{verbatim}
def f():
    try:
        res = int("c6")
    except Exception: return 1
    except ZeroDivisionError: return 7
    else: return 32
    finally: return 20
    return res

print(f())
\end{verbatim}


\vspace{5mm}
               
{\begin{center}Завдання {3}\end{center}}
\vspace{-6mm}

\begin{verbatim}
m=6
while m<=9:
    m+=2
print(m)
\end{verbatim}


\end{minipage}
\vline \hspace{6mm}
\begin{minipage}[t]{10.7cm}
               
{\begin{center}Завдання {4}\end{center}}
\vspace{-6mm}

\begin{verbatim}
for f in range(-5,-5+4,-2):
    if f<=6:
        continue
    print(f, end=' ')
    f=7
else:
    print(13, end=' ')
print(f, end=' ')
\end{verbatim}


\vspace{5mm}
               
{\begin{center}Завдання {5}\end{center}}
\vspace{-6mm}

\begin{verbatim}
def f(n):
    if n > 9:
        f(n // 10)
    else:
        print(n % 10, end=' ')

f(123456)
\end{verbatim}


\end{minipage}

\newpage

{{22.11.2024 Контрольна робота \No 2, варіант  \No \textbf { 18 }\par}}
{
%\begin {large}

Тестова частина. Завдання 1-5: Що буде виведено за виконання (повідомлення інтерпретатора про помилку позначати error)?

Задачі (на звороті): написати функцію для розв'язання задачі.

\vspace{5mm} \begin{minipage}[t]{8.5cm}
               
{\begin{center}Завдання {1}\end{center}}
\vspace{-6mm}

\begin{verbatim}
a,b,c=9,8,3
def h(a):
    global c
    a=5
    b+=1
    c=3
    return a+b+c

a,b,c=1,5,0
print(h(b),a,b,c)
\end{verbatim}


\vspace{5mm}
               
{\begin{center}Завдання {2}\end{center}}
\vspace{-6mm}

\begin{verbatim}
def f():
    try:
        res = int("1")
        return 41
    except BaseException: return 6
    except TypeError: return 9
    return res

print(f())
\end{verbatim}


\vspace{5mm}
               
{\begin{center}Завдання {3}\end{center}}
\vspace{-6mm}

\begin{verbatim}
t=5
while t<11:
    t+=1
print(t)
\end{verbatim}


\end{minipage}
\vline \hspace{6mm}
\begin{minipage}[t]{10.7cm}
               
{\begin{center}Завдання {4}\end{center}}
\vspace{-6mm}

\begin{verbatim}
for a in range(-6,-6+3,2):
    if a>4:
        pass
    print(a, end=' ')
    a=3
else:
    print(a, end=' ')
print(a, end=' ')
\end{verbatim}


\vspace{5mm}
               
{\begin{center}Завдання {5}\end{center}}
\vspace{-6mm}

\begin{verbatim}
def f(n):
    print(n % 10, end=' ')
    if n > 9:
        f(n // 100)
    

f(987654)
\end{verbatim}


\end{minipage}

\newpage

{{22.11.2024 Контрольна робота \No 2, варіант  \No \textbf { 19 }\par}}
{
%\begin {large}

Тестова частина. Завдання 1-5: Що буде виведено за виконання (повідомлення інтерпретатора про помилку позначати error)?

Задачі (на звороті): написати функцію для розв'язання задачі.

\vspace{5mm} \begin{minipage}[t]{8.5cm}
               
{\begin{center}Завдання {1}\end{center}}
\vspace{-6mm}

\begin{verbatim}
a,b,c=6,7,3
def g(b):
    a*=2
    b=3
    c=5
    return a+b+c

a,b,c=7,4,2
print(g(a),a,b,c)
\end{verbatim}


\vspace{5mm}
               
{\begin{center}Завдання {2}\end{center}}
\vspace{-6mm}

\begin{verbatim}
def f():
    try:
        res = int(5//1)
        return 40
    except BaseException: return 4
    except KeyboardInterrupt: return 0
    else: return 35
    return res

print(f())
\end{verbatim}


\vspace{5mm}
               
{\begin{center}Завдання {3}\end{center}}
\vspace{-6mm}

\begin{verbatim}
n=9
while n<=11:
    n+=2
print(n)
\end{verbatim}


\end{minipage}
\vline \hspace{6mm}
\begin{minipage}[t]{10.7cm}
               
{\begin{center}Завдання {4}\end{center}}
\vspace{-6mm}

\begin{verbatim}
for c in range(4,4+5,3):
    if c<8:
        break
    print(c, end=' ')
    c=10
else:
    print(c, end=' ')
print(c, end=' ')
\end{verbatim}


\vspace{5mm}
               
{\begin{center}Завдання {5}\end{center}}
\vspace{-6mm}

\begin{verbatim}
def f(n):
    if n > 9:
        f(n // 100)
    else:
        print(n % 10, end=' ')

f(123456)
\end{verbatim}


\end{minipage}

\newpage

{{22.11.2024 Контрольна робота \No 2, варіант  \No \textbf { 20 }\par}}
{
%\begin {large}

Тестова частина. Завдання 1-5: Що буде виведено за виконання (повідомлення інтерпретатора про помилку позначати error)?

Задачі (на звороті): написати функцію для розв'язання задачі.

\vspace{5mm} \begin{minipage}[t]{8.5cm}
               
{\begin{center}Завдання {1}\end{center}}
\vspace{-6mm}

\begin{verbatim}
a,b,c=9,8,0
def f(b):
    global c
    a=2
    b-=1
    c=4
    return a+b+c

a,b,c=5,6,1
print(f(a),a,b,c)
\end{verbatim}


\vspace{5mm}
               
{\begin{center}Завдання {2}\end{center}}
\vspace{-6mm}

\begin{verbatim}
def f():
    try:
        res = int(8/0)
    except ZeroDivisionError: return 3
    except ValueError: return 2
    finally: return 23
    return res

print(f())
\end{verbatim}


\vspace{5mm}
               
{\begin{center}Завдання {3}\end{center}}
\vspace{-6mm}

\begin{verbatim}
i=8
while i<=5:
    i+=2
print(i)
\end{verbatim}


\end{minipage}
\vline \hspace{6mm}
\begin{minipage}[t]{10.7cm}
               
{\begin{center}Завдання {4}\end{center}}
\vspace{-6mm}

\begin{verbatim}
for d in range(-9,-9+6,-3):
    if d>=6:
        continue
    print(d, end=' ')
    d=-8
else:
    print(13, end=' ')
print(d, end=' ')
\end{verbatim}


\vspace{5mm}
               
{\begin{center}Завдання {5}\end{center}}
\vspace{-6mm}

\begin{verbatim}
def f(n):
    if n <= 9:
        print(n % 10, end=' ')
    else:
        f(n // 10)

f(123456)
\end{verbatim}


\end{minipage}

\newpage

{{22.11.2024 Контрольна робота \No 2, варіант  \No \textbf { 21 }\par}}
{
%\begin {large}

Тестова частина. Завдання 1-5: Що буде виведено за виконання (повідомлення інтерпретатора про помилку позначати error)?

Задачі (на звороті): написати функцію для розв'язання задачі.

\vspace{5mm} \begin{minipage}[t]{8.5cm}
               
{\begin{center}Завдання {1}\end{center}}
\vspace{-6mm}

\begin{verbatim}
a,b,c=8,4,1
def f(a):
    a=5
    b*=1
    c=3
    return a+b+c

a,b,c=3,0,5
print(f(b),a,b,c)
\end{verbatim}


\vspace{5mm}
               
{\begin{center}Завдання {2}\end{center}}
\vspace{-6mm}

\begin{verbatim}
def f():
    try:
        res = int(4%1)
    except Exception: return 8
    except TypeError: return 2
    else: return 31
    finally: return 22
    return res

print(f())
\end{verbatim}


\vspace{5mm}
               
{\begin{center}Завдання {3}\end{center}}
\vspace{-6mm}

\begin{verbatim}
j=0
while j<7:
    j+=1
print(j)
\end{verbatim}


\end{minipage}
\vline \hspace{6mm}
\begin{minipage}[t]{10.7cm}
               
{\begin{center}Завдання {4}\end{center}}
\vspace{-6mm}

\begin{verbatim}
for e in range(-3,-3+6,2):
    if e>=3:
        continue
    print(e, end=' ')
    e=8
else:
    print(e, end=' ')
print(e, end=' ')
\end{verbatim}


\vspace{5mm}
               
{\begin{center}Завдання {5}\end{center}}
\vspace{-6mm}

\begin{verbatim}
def f(n):
    print(n % 10, end=' ')
    if n > 9:
        f(n // 10)
    

f(543216)
\end{verbatim}


\end{minipage}

\newpage

{{22.11.2024 Контрольна робота \No 2, варіант  \No \textbf { 22 }\par}}
{
%\begin {large}

Тестова частина. Завдання 1-5: Що буде виведено за виконання (повідомлення інтерпретатора про помилку позначати error)?

Задачі (на звороті): написати функцію для розв'язання задачі.

\vspace{5mm} \begin{minipage}[t]{8.5cm}
               
{\begin{center}Завдання {1}\end{center}}
\vspace{-6mm}

\begin{verbatim}
a,b,c=2,3,1
def h(b):
    a=3
    b=5
    c=2
    return a+b+c

a,b,c=8,7,4
print(h(a),a,b,c)
\end{verbatim}


\vspace{5mm}
               
{\begin{center}Завдання {2}\end{center}}
\vspace{-6mm}

\begin{verbatim}
def f():
    try:
        res = 6!=8
        return 43
    except ZeroDivisionError: return 7
    except KeyboardInterrupt: return 3
    return res

print(f())
\end{verbatim}


\vspace{5mm}
               
{\begin{center}Завдання {3}\end{center}}
\vspace{-6mm}

\begin{verbatim}
n=10
while n>=8:
    n+=-3
print(n)
\end{verbatim}


\end{minipage}
\vline \hspace{6mm}
\begin{minipage}[t]{10.7cm}
               
{\begin{center}Завдання {4}\end{center}}
\vspace{-6mm}

\begin{verbatim}
for b in range(-1,-1+5,3):
    if b>4:
        continue
    print(b, end=' ')
    b=-2
else:
    print(b, end=' ')
print(b, end=' ')
\end{verbatim}


\vspace{5mm}
               
{\begin{center}Завдання {5}\end{center}}
\vspace{-6mm}

\begin{verbatim}
def f(n):
    if n > 9:
        f(n // 10)
    print(n % 10, end=' ')
    

f(123456)
\end{verbatim}


\end{minipage}

\newpage

{{22.11.2024 Контрольна робота \No 2, варіант  \No \textbf { 23 }\par}}
{
%\begin {large}

Тестова частина. Завдання 1-5: Що буде виведено за виконання (повідомлення інтерпретатора про помилку позначати error)?

Задачі (на звороті): написати функцію для розв'язання задачі.

\vspace{5mm} \begin{minipage}[t]{8.5cm}
               
{\begin{center}Завдання {1}\end{center}}
\vspace{-6mm}

\begin{verbatim}
a,b,c=0,6,5
def h(a):
    global c
    a=4
    b=1
    c=4
    return a+b+c

a,b,c=9,3,8
print(h(a),a,b,c)
\end{verbatim}


\vspace{5mm}
               
{\begin{center}Завдання {2}\end{center}}
\vspace{-6mm}

\begin{verbatim}
def f():
    try:
        res = int(9%0)
        return 45
    except Exception: return 5
    except BaseException: return 1
    finally: return 25
    return res

print(f())
\end{verbatim}


\vspace{5mm}
               
{\begin{center}Завдання {3}\end{center}}
\vspace{-6mm}

\begin{verbatim}
k=5
while k>6:
    k+=-3
print(k)
\end{verbatim}


\end{minipage}
\vline \hspace{6mm}
\begin{minipage}[t]{10.7cm}
               
{\begin{center}Завдання {4}\end{center}}
\vspace{-6mm}

\begin{verbatim}
for a in range(1,1+2,-3):
    if a<=7:
        pass
    print(a, end=' ')
    a=-5
    if a>=3:
        break
else:
    print(a, end=' ')
print(a, end=' ')
\end{verbatim}


\vspace{5mm}
               
{\begin{center}Завдання {5}\end{center}}
\vspace{-6mm}

\begin{verbatim}
def f(n):
    if n <= 9:
        print(n % 10, end=' ')
    else:
        f(n // 10)

f(123456)
\end{verbatim}


\end{minipage}

\newpage

{{22.11.2024 Контрольна робота \No 2, варіант  \No \textbf { 24 }\par}}
{
%\begin {large}

Тестова частина. Завдання 1-5: Що буде виведено за виконання (повідомлення інтерпретатора про помилку позначати error)?

Задачі (на звороті): написати функцію для розв'язання задачі.

\vspace{5mm} \begin{minipage}[t]{8.5cm}
               
{\begin{center}Завдання {1}\end{center}}
\vspace{-6mm}

\begin{verbatim}
a,b,c=2,9,6
def h(b):
    global c
    a-=2
    b=4
    c=1
    return a+b+c

a,b,c=7,3,2
print(h(b),a,b,c)
\end{verbatim}


\vspace{5mm}
               
{\begin{center}Завдання {2}\end{center}}
\vspace{-6mm}

\begin{verbatim}
def f():
    try:
        res = int("0")
    except ValueError: return 0
    except TypeError: return 7
    else: return 30
    return res

print(f())
\end{verbatim}


\vspace{5mm}
               
{\begin{center}Завдання {3}\end{center}}
\vspace{-6mm}

\begin{verbatim}
t=8
while t>=5:
    t+=-2
print(t)
\end{verbatim}


\end{minipage}
\vline \hspace{6mm}
\begin{minipage}[t]{10.7cm}
               
{\begin{center}Завдання {4}\end{center}}
\vspace{-6mm}

\begin{verbatim}
for f in range(-4,-4+3,-2):
    if f<5:
        break
    print(f, end=' ')
    f=10
else:
    print(f, end=' ')
print(f, end=' ')
\end{verbatim}


\vspace{5mm}
               
{\begin{center}Завдання {5}\end{center}}
\vspace{-6mm}

\begin{verbatim}
def f(n):
    print(n % 2, end=' ')
    if n > 2:
        f(n // 2)
    

f(16)
\end{verbatim}


\end{minipage}

\newpage

{{22.11.2024 Контрольна робота \No 2, варіант  \No \textbf { 25 }\par}}
{
%\begin {large}

Тестова частина. Завдання 1-5: Що буде виведено за виконання (повідомлення інтерпретатора про помилку позначати error)?

Задачі (на звороті): написати функцію для розв'язання задачі.

\vspace{5mm} \begin{minipage}[t]{8.5cm}
               
{\begin{center}Завдання {1}\end{center}}
\vspace{-6mm}

\begin{verbatim}
a,b,c=8,5,0
def g(a):
    a=5
    b+=4
    c=3
    return a+b+c

a,b,c=6,7,4
print(g(a),a,b,c)
\end{verbatim}


\vspace{5mm}
               
{\begin{center}Завдання {2}\end{center}}
\vspace{-6mm}

\begin{verbatim}
def f():
    try:
        res = int("b3")
        return 45
    except KeyboardInterrupt: return 1
    except Exception: return 6
    return res

print(f())
\end{verbatim}


\vspace{5mm}
               
{\begin{center}Завдання {3}\end{center}}
\vspace{-6mm}

\begin{verbatim}
j=6
while j>=1:
    j+=-3
print(j)
\end{verbatim}


\end{minipage}
\vline \hspace{6mm}
\begin{minipage}[t]{10.7cm}
               
{\begin{center}Завдання {4}\end{center}}
\vspace{-6mm}

\begin{verbatim}
for a in range(-2,-2+4,-3):
    if a<6:
        break
    print(a, end=' ')
    a=-4
else:
    print(a, end=' ')
print(a, end=' ')
\end{verbatim}


\vspace{5mm}
               
{\begin{center}Завдання {5}\end{center}}
\vspace{-6mm}

\begin{verbatim}
def f(n):
    if n > 9:
        return f(n // 10) + 1
    else:
        return 1

print(f(123456))
\end{verbatim}


\end{minipage}

\newpage

{{22.11.2024 Контрольна робота \No 2, варіант  \No \textbf { 26 }\par}}
{
%\begin {large}

Тестова частина. Завдання 1-5: Що буде виведено за виконання (повідомлення інтерпретатора про помилку позначати error)?

Задачі (на звороті): написати функцію для розв'язання задачі.

\vspace{5mm} \begin{minipage}[t]{8.5cm}
               
{\begin{center}Завдання {1}\end{center}}
\vspace{-6mm}

\begin{verbatim}
a,b,c=9,1,8
def h(a):
    a=2
    b*=1
    c=2
    return a+b+c

a,b,c=3,1,9
print(h(a),a,b,c)
\end{verbatim}


\vspace{5mm}
               
{\begin{center}Завдання {2}\end{center}}
\vspace{-6mm}

\begin{verbatim}
def f():
    try:
        res = int("d5")
    except ValueError: return 4
    except TypeError: return 8
    else: return 30
    finally: return 23
    return res

print(f())
\end{verbatim}


\vspace{5mm}
               
{\begin{center}Завдання {3}\end{center}}
\vspace{-6mm}

\begin{verbatim}
m=2
while m<8:
    m+=2
print(m)
\end{verbatim}


\end{minipage}
\vline \hspace{6mm}
\begin{minipage}[t]{10.7cm}
               
{\begin{center}Завдання {4}\end{center}}
\vspace{-6mm}

\begin{verbatim}
for f in range(-5,-5+2,3):
    if f>4:
        continue
    print(f, end=' ')
    f=-10
else:
    print(13, end=' ')
print(f, end=' ')
\end{verbatim}


\vspace{5mm}
               
{\begin{center}Завдання {5}\end{center}}
\vspace{-6mm}

\begin{verbatim}
def f(n):
    if n <= 9:
        print(n % 10, end=' ')
    else:
        f(n // 10)

f(543216)
\end{verbatim}


\end{minipage}

\newpage

{{22.11.2024 Контрольна робота \No 2, варіант  \No \textbf { 27 }\par}}
{
%\begin {large}

Тестова частина. Завдання 1-5: Що буде виведено за виконання (повідомлення інтерпретатора про помилку позначати error)?

Задачі (на звороті): написати функцію для розв'язання задачі.

\vspace{5mm} \begin{minipage}[t]{8.5cm}
               
{\begin{center}Завдання {1}\end{center}}
\vspace{-6mm}

\begin{verbatim}
a,b,c=7,0,5
def g(b):
    global c
    a-=5
    b=3
    c=4
    return a+b+c

a,b,c=6,2,4
print(g(a),a,b,c)
\end{verbatim}


\vspace{5mm}
               
{\begin{center}Завдання {2}\end{center}}
\vspace{-6mm}

\begin{verbatim}
def f():
    try:
        res = 2>9
        return 41
    except ZeroDivisionError: return 9
    except BaseException: return 2
    return res

print(f())
\end{verbatim}


\vspace{5mm}
               
{\begin{center}Завдання {3}\end{center}}
\vspace{-6mm}

\begin{verbatim}
m=10
while m>0:
    m+=-2
print(m)
\end{verbatim}


\end{minipage}
\vline \hspace{6mm}
\begin{minipage}[t]{10.7cm}
               
{\begin{center}Завдання {4}\end{center}}
\vspace{-6mm}

\begin{verbatim}
for e in range(6,6+5,-2):
    if e<=8:
        continue
    print(e, end=' ')
    e=-8
else:
    print(e, end=' ')
print(e, end=' ')
\end{verbatim}


\vspace{5mm}
               
{\begin{center}Завдання {5}\end{center}}
\vspace{-6mm}

\begin{verbatim}
def f(n):
    if n > 2:
        f(n // 2)
    print(n % 2, end=' ')
    

f(16)
\end{verbatim}


\end{minipage}

\newpage

{{22.11.2024 Контрольна робота \No 2, варіант  \No \textbf { 28 }\par}}
{
%\begin {large}

Тестова частина. Завдання 1-5: Що буде виведено за виконання (повідомлення інтерпретатора про помилку позначати error)?

Задачі (на звороті): написати функцію для розв'язання задачі.

\vspace{5mm} \begin{minipage}[t]{8.5cm}
               
{\begin{center}Завдання {1}\end{center}}
\vspace{-6mm}

\begin{verbatim}
a,b,c=5,2,6
def f(b):
    a+=4
    b=5
    c=2
    return a+b+c

a,b,c=8,0,9
print(f(b),a,b,c)
\end{verbatim}


\vspace{5mm}
               
{\begin{center}Завдання {2}\end{center}}
\vspace{-6mm}

\begin{verbatim}
def f():
    try:
        res = int(7//2)
    except ValueError: return 5
    except TypeError: return 4
    else: return 34
    finally: return 20
    return res

print(f())
\end{verbatim}


\vspace{5mm}
               
{\begin{center}Завдання {3}\end{center}}
\vspace{-6mm}

\begin{verbatim}
k=5
while k<=12:
    k+=1
print(k)
\end{verbatim}


\end{minipage}
\vline \hspace{6mm}
\begin{minipage}[t]{10.7cm}
               
{\begin{center}Завдання {4}\end{center}}
\vspace{-6mm}

\begin{verbatim}
for d in range(4,4+3,2):
    if d>=3:
        pass
    print(d, end=' ')
    d=1
else:
    print(d, end=' ')
print(d, end=' ')
\end{verbatim}


\vspace{5mm}
               
{\begin{center}Завдання {5}\end{center}}
\vspace{-6mm}

\begin{verbatim}
def f(n):
    if n > 2:
        f(n // 2)
    else:
        print(n % 2, end=' ')

f(16)
\end{verbatim}


\end{minipage}

\newpage

{{22.11.2024 Контрольна робота \No 2, варіант  \No \textbf { 29 }\par}}
{
%\begin {large}

Тестова частина. Завдання 1-5: Що буде виведено за виконання (повідомлення інтерпретатора про помилку позначати error)?

Задачі (на звороті): написати функцію для розв'язання задачі.

\vspace{5mm} \begin{minipage}[t]{8.5cm}
               
{\begin{center}Завдання {1}\end{center}}
\vspace{-6mm}

\begin{verbatim}
a,b,c=0,4,2
def g(b):
    a=3
    b=1
    c=5
    return a+b+c

a,b,c=7,5,1
print(g(b),a,b,c)
\end{verbatim}


\vspace{5mm}
               
{\begin{center}Завдання {2}\end{center}}
\vspace{-6mm}

\begin{verbatim}
def f():
    try:
        res = int("3")
        return 43
    except KeyboardInterrupt: return 9
    except BaseException: return 0
    else: return 35
    finally: return 21
    return res

print(f())
\end{verbatim}


\vspace{5mm}
               
{\begin{center}Завдання {3}\end{center}}
\vspace{-6mm}

\begin{verbatim}
t=9
while t>7:
    t+=-3
print(t)
\end{verbatim}


\end{minipage}
\vline \hspace{6mm}
\begin{minipage}[t]{10.7cm}
               
{\begin{center}Завдання {4}\end{center}}
\vspace{-6mm}

\begin{verbatim}
for c in range(9,9+6,2):
    if c>=7:
        continue
    print(c, end=' ')
    c=5
else:
    print(c, end=' ')
print(c, end=' ')
\end{verbatim}


\vspace{5mm}
               
{\begin{center}Завдання {5}\end{center}}
\vspace{-6mm}

\begin{verbatim}
def f(n):
    if n > 9:
            f(n // 100)
    print(n % 10, end=' ')
    

f(987654)
\end{verbatim}


\end{minipage}

\newpage

{{22.11.2024 Контрольна робота \No 2, варіант  \No \textbf { 30 }\par}}
{
%\begin {large}

Тестова частина. Завдання 1-5: Що буде виведено за виконання (повідомлення інтерпретатора про помилку позначати error)?

Задачі (на звороті): написати функцію для розв'язання задачі.

\vspace{5mm} \begin{minipage}[t]{8.5cm}
               
{\begin{center}Завдання {1}\end{center}}
\vspace{-6mm}

\begin{verbatim}
a,b,c=7,1,3
def h(a):
    global c
    a=3
    b*=1
    c=4
    return a+b+c

a,b,c=4,8,3
print(h(b),a,b,c)
\end{verbatim}


\vspace{5mm}
               
{\begin{center}Завдання {2}\end{center}}
\vspace{-6mm}

\begin{verbatim}
def f():
    try:
        res = int(8/0.0)
    except ZeroDivisionError: return 6
    except Exception: return 1
    return res

print(f())
\end{verbatim}


\vspace{5mm}
               
{\begin{center}Завдання {3}\end{center}}
\vspace{-6mm}

\begin{verbatim}
j=6
while j<9:
    j+=2
print(j)
\end{verbatim}


\end{minipage}
\vline \hspace{6mm}
\begin{minipage}[t]{10.7cm}
               
{\begin{center}Завдання {4}\end{center}}
\vspace{-6mm}

\begin{verbatim}
for b in range(10,10+4,3):
    if b<5:
        break
    print(b, end=' ')
    b=8
else:
    print(b, end=' ')
print(b, end=' ')
\end{verbatim}


\vspace{5mm}
               
{\begin{center}Завдання {5}\end{center}}
\vspace{-6mm}

\begin{verbatim}
def f(n):
    if n > 9:
        f(n // 100)
    print(n % 10, end=' ')
    

f(123456)
\end{verbatim}


\end{minipage}

\newpage

{{22.11.2024 Контрольна робота \No 2, варіант  \No \textbf { 31 }\par}}
{
%\begin {large}

Тестова частина. Завдання 1-5: Що буде виведено за виконання (повідомлення інтерпретатора про помилку позначати error)?

Задачі (на звороті): написати функцію для розв'язання задачі.

\vspace{5mm} \begin{minipage}[t]{8.5cm}
               
{\begin{center}Завдання {1}\end{center}}
\vspace{-6mm}

\begin{verbatim}
a,b,c=6,9,4
def f(b):
    a=2
    b=1
    c=4
    return a+b+c

a,b,c=0,8,7
print(f(b),a,b,c)
\end{verbatim}


\vspace{5mm}
               
{\begin{center}Завдання {2}\end{center}}
\vspace{-6mm}

\begin{verbatim}
def f():
    try:
        res = 9<=6
    except ZeroDivisionError: return 0
    except Exception: return 8
    else: return 33
    finally: return 24
    return res

print(f())
\end{verbatim}


\vspace{5mm}
               
{\begin{center}Завдання {3}\end{center}}
\vspace{-6mm}

\begin{verbatim}
i=4
while i<=6:
    i+=1
print(i)
\end{verbatim}


\end{minipage}
\vline \hspace{6mm}
\begin{minipage}[t]{10.7cm}
               
{\begin{center}Завдання {4}\end{center}}
\vspace{-6mm}

\begin{verbatim}
for e in range(-10,-10+6,-2):
    if e<=4:
        pass
    print(e, end=' ')
    e=9
else:
    print(e, end=' ')
print(e, end=' ')
\end{verbatim}


\vspace{5mm}
               
{\begin{center}Завдання {5}\end{center}}
\vspace{-6mm}

\begin{verbatim}
def f(n):
    print(n % 10, end=' ')
    if n > 9:
        f(n // 10)
    

f(123456)
\end{verbatim}


\end{minipage}

\newpage

{{22.11.2024 Контрольна робота \No 2, варіант  \No \textbf { 32 }\par}}
{
%\begin {large}

Тестова частина. Завдання 1-5: Що буде виведено за виконання (повідомлення інтерпретатора про помилку позначати error)?

Задачі (на звороті): написати функцію для розв'язання задачі.

\vspace{5mm} \begin{minipage}[t]{8.5cm}
               
{\begin{center}Завдання {1}\end{center}}
\vspace{-6mm}

\begin{verbatim}
a,b,c=1,6,2
def h(a):
    global c
    a=3
    b=5
    c=2
    return a+b+c

a,b,c=3,9,5
print(h(a),a,b,c)
\end{verbatim}


\vspace{5mm}
               
{\begin{center}Завдання {2}\end{center}}
\vspace{-6mm}

\begin{verbatim}
def f():
    try:
        res = int("2")
        return 40
    except ValueError: return 6
    except TypeError: return 7
    return res

print(f())
\end{verbatim}


\vspace{5mm}
               
{\begin{center}Завдання {3}\end{center}}
\vspace{-6mm}

\begin{verbatim}
n=7
while n<=10:
    n+=1
print(n)
\end{verbatim}


\end{minipage}
\vline \hspace{6mm}
\begin{minipage}[t]{10.7cm}
               
{\begin{center}Завдання {4}\end{center}}
\vspace{-6mm}

\begin{verbatim}
for a in range(0,0+2,-3):
    if a>8:
        continue
    print(a, end=' ')
    a=-3
    if a<5:
        break
else:
    print(a, end=' ')
print(a, end=' ')
\end{verbatim}


\vspace{5mm}
               
{\begin{center}Завдання {5}\end{center}}
\vspace{-6mm}

\begin{verbatim}
def f(n):
    if n <= 9:
        print(n % 10, end=' ')
    else:
        f(n // 10)

f(987651)
\end{verbatim}


\end{minipage}

\newpage

{{22.11.2024 Контрольна робота \No 2, варіант  \No \textbf { 33 }\par}}
{
%\begin {large}

Тестова частина. Завдання 1-5: Що буде виведено за виконання (повідомлення інтерпретатора про помилку позначати error)?

Задачі (на звороті): написати функцію для розв'язання задачі.

\vspace{5mm} \begin{minipage}[t]{8.5cm}
               
{\begin{center}Завдання {1}\end{center}}
\vspace{-6mm}

\begin{verbatim}
a,b,c=4,1,0
def f(b):
    global c
    a-=2
    b=5
    c=3
    return a+b+c

a,b,c=7,6,5
print(f(a),a,b,c)
\end{verbatim}


\vspace{5mm}
               
{\begin{center}Завдання {2}\end{center}}
\vspace{-6mm}

\begin{verbatim}
def f():
    try:
        res = int(5%2)
    except BaseException: return 1
    except KeyboardInterrupt: return 3
    else: return 31
    return res

print(f())
\end{verbatim}


\vspace{5mm}
               
{\begin{center}Завдання {3}\end{center}}
\vspace{-6mm}

\begin{verbatim}
n=14
while n>=4:
    n+=-2
print(n)
\end{verbatim}


\end{minipage}
\vline \hspace{6mm}
\begin{minipage}[t]{10.7cm}
               
{\begin{center}Завдання {4}\end{center}}
\vspace{-6mm}

\begin{verbatim}
for c in range(-7,-7+5,-3):
    if c>5:
        continue
    print(c, end=' ')
    c=-6
else:
    print(13, end=' ')
print(c, end=' ')
\end{verbatim}


\vspace{5mm}
               
{\begin{center}Завдання {5}\end{center}}
\vspace{-6mm}

\begin{verbatim}
def f(n):
    if n > 9:
        f(n // 10)
    else:
        print(n % 10, end=' ')

f(543216)
\end{verbatim}


\end{minipage}

\newpage

{{22.11.2024 Контрольна робота \No 2, варіант  \No \textbf { 34 }\par}}
{
%\begin {large}

Тестова частина. Завдання 1-5: Що буде виведено за виконання (повідомлення інтерпретатора про помилку позначати error)?

Задачі (на звороті): написати функцію для розв'язання задачі.

\vspace{5mm} \begin{minipage}[t]{8.5cm}
               
{\begin{center}Завдання {1}\end{center}}
\vspace{-6mm}

\begin{verbatim}
a,b,c=2,9,2
def g(a):
    a+=1
    b=3
    c=4
    return a+b+c

a,b,c=0,4,3
print(g(b),a,b,c)
\end{verbatim}


\vspace{5mm}
               
{\begin{center}Завдання {2}\end{center}}
\vspace{-6mm}

\begin{verbatim}
def f():
    try:
        res = int(4/0)
        return 42
    except ZeroDivisionError: return 5
    except ValueError: return 6
    finally: return 25
    return res

print(f())
\end{verbatim}


\vspace{5mm}
               
{\begin{center}Завдання {3}\end{center}}
\vspace{-6mm}

\begin{verbatim}
i=15
while i>=2:
    i+=-2
print(i)
\end{verbatim}


\end{minipage}
\vline \hspace{6mm}
\begin{minipage}[t]{10.7cm}
               
{\begin{center}Завдання {4}\end{center}}
\vspace{-6mm}

\begin{verbatim}
for f in range(-9,-9+3,2):
    if f>=3:
        break
    print(f, end=' ')
    f=4
else:
    print(f, end=' ')
print(f, end=' ')
\end{verbatim}


\vspace{5mm}
               
{\begin{center}Завдання {5}\end{center}}
\vspace{-6mm}

\begin{verbatim}
def f(n):
    print(n % 10, end=' ')
    if n > 9:
        f(n // 100)
    

f(987654)
\end{verbatim}


\end{minipage}

\newpage

{{22.11.2024 Контрольна робота \No 2, варіант  \No \textbf { 35 }\par}}
{
%\begin {large}

Тестова частина. Завдання 1-5: Що буде виведено за виконання (повідомлення інтерпретатора про помилку позначати error)?

Задачі (на звороті): написати функцію для розв'язання задачі.

\vspace{5mm} \begin{minipage}[t]{8.5cm}
               
{\begin{center}Завдання {1}\end{center}}
\vspace{-6mm}

\begin{verbatim}
a,b,c=1,8,5
def f(a):
    global c
    a+=1
    b=5
    c=2
    return a+b+c

a,b,c=6,7,9
print(f(a),a,b,c)
\end{verbatim}


\vspace{5mm}
               
{\begin{center}Завдання {2}\end{center}}
\vspace{-6mm}

\begin{verbatim}
def f():
    try:
        res = int("a0")
    except Exception: return 8
    except TypeError: return 2
    else: return 32
    finally: return 22
    return res

print(f())
\end{verbatim}


\vspace{5mm}
               
{\begin{center}Завдання {3}\end{center}}
\vspace{-6mm}

\begin{verbatim}
k=13
while k>8:
    k+=-3
print(k)
\end{verbatim}


\end{minipage}
\vline \hspace{6mm}
\begin{minipage}[t]{10.7cm}
               
{\begin{center}Завдання {4}\end{center}}
\vspace{-6mm}

\begin{verbatim}
for d in range(8,8+4,-2):
    if d<7:
        pass
    print(d, end=' ')
    d=-7
else:
    print(d, end=' ')
print(d, end=' ')
\end{verbatim}


\vspace{5mm}
               
{\begin{center}Завдання {5}\end{center}}
\vspace{-6mm}

\begin{verbatim}
def f(n):
    if n > 9:
        f(n // 100)
    else:
        print(n % 10, end=' ')

f(123456)
\end{verbatim}


\end{minipage}

\newpage

{{22.11.2024 Контрольна робота \No 2, варіант  \No \textbf { 36 }\par}}
{
%\begin {large}

Тестова частина. Завдання 1-5: Що буде виведено за виконання (повідомлення інтерпретатора про помилку позначати error)?

Задачі (на звороті): написати функцію для розв'язання задачі.

\vspace{5mm} \begin{minipage}[t]{8.5cm}
               
{\begin{center}Завдання {1}\end{center}}
\vspace{-6mm}

\begin{verbatim}
a,b,c=9,4,1
def h(b):
    a=3
    b*=1
    c=2
    return a+b+c

a,b,c=5,2,6
print(h(b),a,b,c)
\end{verbatim}


\vspace{5mm}
               
{\begin{center}Завдання {2}\end{center}}
\vspace{-6mm}

\begin{verbatim}
def f():
    try:
        res = int("c3")
        return 44
    except KeyboardInterrupt: return 7
    except BaseException: return 1
    return res

print(f())
\end{verbatim}


\vspace{5mm}
               
{\begin{center}Завдання {3}\end{center}}
\vspace{-6mm}

\begin{verbatim}
k=11
while k>3:
    k+=-3
print(k)
\end{verbatim}


\end{minipage}
\vline \hspace{6mm}
\begin{minipage}[t]{10.7cm}
               
{\begin{center}Завдання {4}\end{center}}
\vspace{-6mm}

\begin{verbatim}
for b in range(-8,-8+5,3):
    if b<=6:
        continue
    print(b, end=' ')
    b=-9
else:
    print(b, end=' ')
print(b, end=' ')
\end{verbatim}


\vspace{5mm}
               
{\begin{center}Завдання {5}\end{center}}
\vspace{-6mm}

\begin{verbatim}
def f(n):
    print(n % 10, end=' ')
    if n > 9:
           f(n // 100)
    

f(123456)
\end{verbatim}


\end{minipage}

\newpage

{{22.11.2024 Контрольна робота \No 2, варіант  \No \textbf { 37 }\par}}
{
%\begin {large}

Тестова частина. Завдання 1-5: Що буде виведено за виконання (повідомлення інтерпретатора про помилку позначати error)?

Задачі (на звороті): написати функцію для розв'язання задачі.

\vspace{5mm} \begin{minipage}[t]{8.5cm}
               
{\begin{center}Завдання {1}\end{center}}
\vspace{-6mm}

\begin{verbatim}
a,b,c=0,7,3
def g(b):
    a=5
    b-=4
    c=1
    return a+b+c

a,b,c=8,7,4
print(g(a),a,b,c)
\end{verbatim}


\vspace{5mm}
               
{\begin{center}Завдання {2}\end{center}}
\vspace{-6mm}

\begin{verbatim}
def f():
    try:
        res = int("9")
        return 42
    except KeyboardInterrupt: return 4
    except ValueError: return 2
    else: return 35
    finally: return 22
    return res

print(f())
\end{verbatim}


\vspace{5mm}
               
{\begin{center}Завдання {3}\end{center}}
\vspace{-6mm}

\begin{verbatim}
t=1
while t<13:
    t+=2
print(t)
\end{verbatim}


\end{minipage}
\vline \hspace{6mm}
\begin{minipage}[t]{10.7cm}
               
{\begin{center}Завдання {4}\end{center}}
\vspace{-6mm}

\begin{verbatim}
for b in range(2,2+3,2):
    if b<6:
        continue
    print(b, end=' ')
    b=7
    if b<=4:
        break
else:
    print(b, end=' ')
print(b, end=' ')
\end{verbatim}


\vspace{5mm}
               
{\begin{center}Завдання {5}\end{center}}
\vspace{-6mm}

\begin{verbatim}
def f(n):
    if n > 9:
        return f(n // 100) + 1
    else:
        return 1

print(f(123456))
\end{verbatim}


\end{minipage}

\newpage

{{22.11.2024 Контрольна робота \No 2, варіант  \No \textbf { 38 }\par}}
{
%\begin {large}

Тестова частина. Завдання 1-5: Що буде виведено за виконання (повідомлення інтерпретатора про помилку позначати error)?

Задачі (на звороті): написати функцію для розв'язання задачі.

\vspace{5mm} \begin{minipage}[t]{8.5cm}
               
{\begin{center}Завдання {1}\end{center}}
\vspace{-6mm}

\begin{verbatim}
a,b,c=6,8,0
def f(a):
    global c
    a*=4
    b=2
    c=5
    return a+b+c

a,b,c=2,3,1
print(f(b),a,b,c)
\end{verbatim}


\vspace{5mm}
               
{\begin{center}Завдання {2}\end{center}}
\vspace{-6mm}

\begin{verbatim}
def f():
    try:
        res = int(0//1)
    except ZeroDivisionError: return 8
    except BaseException: return 6
    return res

print(f())
\end{verbatim}


\vspace{5mm}
               
{\begin{center}Завдання {3}\end{center}}
\vspace{-6mm}

\begin{verbatim}
m=12
while m>=9:
    m+=-2
print(m)
\end{verbatim}


\end{minipage}
\vline \hspace{6mm}
\begin{minipage}[t]{10.7cm}
               
{\begin{center}Завдання {4}\end{center}}
\vspace{-6mm}

\begin{verbatim}
for f in range(-6,-6+2,-2):
    if f>=5:
        pass
    print(f, end=' ')
    f=3
else:
    print(f, end=' ')
print(f, end=' ')
\end{verbatim}


\vspace{5mm}
               
{\begin{center}Завдання {5}\end{center}}
\vspace{-6mm}

\begin{verbatim}
def f(n):
    if n > 9:
        f(n // 10)
    else:
        print(n % 10, end=' ')

f(123456)
\end{verbatim}


\end{minipage}

\newpage

{{22.11.2024 Контрольна робота \No 2, варіант  \No \textbf { 39 }\par}}
{
%\begin {large}

Тестова частина. Завдання 1-5: Що буде виведено за виконання (повідомлення інтерпретатора про помилку позначати error)?

Задачі (на звороті): написати функцію для розв'язання задачі.

\vspace{5mm} \begin{minipage}[t]{8.5cm}
               
{\begin{center}Завдання {1}\end{center}}
\vspace{-6mm}

\begin{verbatim}
a,b,c=9,5,3
def g(b):
    global c
    a=3
    b+=1
    c=3
    return a+b+c

a,b,c=8,6,7
print(g(b),a,b,c)
\end{verbatim}


\vspace{5mm}
               
{\begin{center}Завдання {2}\end{center}}
\vspace{-6mm}

\begin{verbatim}
def f():
    try:
        res = 7>=2
    except TypeError: return 3
    except Exception: return 4
    return res

print(f())
\end{verbatim}


\vspace{5mm}
               
{\begin{center}Завдання {3}\end{center}}
\vspace{-6mm}

\begin{verbatim}
k=3
while k<6:
    k+=2
print(k)
\end{verbatim}


\end{minipage}
\vline \hspace{6mm}
\begin{minipage}[t]{10.7cm}
               
{\begin{center}Завдання {4}\end{center}}
\vspace{-6mm}

\begin{verbatim}
for c in range(3,3+4,-3):
    if c>3:
        break
    print(c, end=' ')
    c=2
else:
    print(13, end=' ')
print(c, end=' ')
\end{verbatim}


\vspace{5mm}
               
{\begin{center}Завдання {5}\end{center}}
\vspace{-6mm}

\begin{verbatim}
def f(n):
    if n > 9:
        f(n // 10)
    print(n % 10, end=' ')
    

f(543216)
\end{verbatim}


\end{minipage}

\newpage

{{22.11.2024 Контрольна робота \No 2, варіант  \No \textbf { 40 }\par}}
{
%\begin {large}

Тестова частина. Завдання 1-5: Що буде виведено за виконання (повідомлення інтерпретатора про помилку позначати error)?

Задачі (на звороті): написати функцію для розв'язання задачі.

\vspace{5mm} \begin{minipage}[t]{8.5cm}
               
{\begin{center}Завдання {1}\end{center}}
\vspace{-6mm}

\begin{verbatim}
a,b,c=4,1,9
def h(a):
    a=2
    b-=5
    c=4
    return a+b+c

a,b,c=5,0,2
print(h(a),a,b,c)
\end{verbatim}


\vspace{5mm}
               
{\begin{center}Завдання {2}\end{center}}
\vspace{-6mm}

\begin{verbatim}
def f():
    try:
        res = int(5/0.0)
        return 45
    except TypeError: return 9
    except ValueError: return 1
    else: return 33
    finally: return 23
    return res

print(f())
\end{verbatim}


\vspace{5mm}
               
{\begin{center}Завдання {3}\end{center}}
\vspace{-6mm}

\begin{verbatim}
i=9
while i<=13:
    i+=1
print(i)
\end{verbatim}


\end{minipage}
\vline \hspace{6mm}
\begin{minipage}[t]{10.7cm}
               
{\begin{center}Завдання {4}\end{center}}
\vspace{-6mm}

\begin{verbatim}
for e in range(-3,-3+6,3):
    if e<=8:
        continue
    print(e, end=' ')
    e=-1
else:
    print(e, end=' ')
print(e, end=' ')
\end{verbatim}


\vspace{5mm}
               
{\begin{center}Завдання {5}\end{center}}
\vspace{-6mm}

\begin{verbatim}
def f(n):
    if n > 9:
        f(n // 100)
    else:
        print(n % 10, end=' ')

f(987654)
\end{verbatim}


\end{minipage}

\newpage

{{22.11.2024 Контрольна робота \No 2, варіант  \No \textbf { 41 }\par}}
{
%\begin {large}

Тестова частина. Завдання 1-5: Що буде виведено за виконання (повідомлення інтерпретатора про помилку позначати error)?

Задачі (на звороті): написати функцію для розв'язання задачі.

\vspace{5mm} \begin{minipage}[t]{8.5cm}
               
{\begin{center}Завдання {1}\end{center}}
\vspace{-6mm}

\begin{verbatim}
a,b,c=6,1,0
def g(b):
    a=2
    b=4
    c=1
    return a+b+c

a,b,c=2,3,5
print(g(b),a,b,c)
\end{verbatim}


\vspace{5mm}
               
{\begin{center}Завдання {2}\end{center}}
\vspace{-6mm}

\begin{verbatim}
def f():
    try:
        res = int(7%0.0)
        return 40
    except Exception: return 9
    except KeyboardInterrupt: return 5
    finally: return 24
    return res

print(f())
\end{verbatim}


\vspace{5mm}
               
{\begin{center}Завдання {3}\end{center}}
\vspace{-6mm}

\begin{verbatim}
t=5
while t>=2:
    t+=-2
print(t)
\end{verbatim}


\end{minipage}
\vline \hspace{6mm}
\begin{minipage}[t]{10.7cm}
               
{\begin{center}Завдання {4}\end{center}}
\vspace{-6mm}

\begin{verbatim}
for d in range(7,7+2,2):
    if d>=7:
        pass
    print(d, end=' ')
    d=6
else:
    print(d, end=' ')
print(d, end=' ')
\end{verbatim}


\vspace{5mm}
               
{\begin{center}Завдання {5}\end{center}}
\vspace{-6mm}

\begin{verbatim}
def f(n):
    print(n % 10, end=' ')
    if n > 9:
           f(n // 100)
    

f(123456)
\end{verbatim}


\end{minipage}

\newpage

{{22.11.2024 Контрольна робота \No 2, варіант  \No \textbf { 42 }\par}}
{
%\begin {large}

Тестова частина. Завдання 1-5: Що буде виведено за виконання (повідомлення інтерпретатора про помилку позначати error)?

Задачі (на звороті): написати функцію для розв'язання задачі.

\vspace{5mm} \begin{minipage}[t]{8.5cm}
               
{\begin{center}Завдання {1}\end{center}}
\vspace{-6mm}

\begin{verbatim}
a,b,c=1,3,7
def g(b):
    a*=1
    b=2
    c=5
    return a+b+c

a,b,c=9,5,8
print(g(b),a,b,c)
\end{verbatim}


\vspace{5mm}
               
{\begin{center}Завдання {2}\end{center}}
\vspace{-6mm}

\begin{verbatim}
def f():
    try:
        res = int(4//1)
    except BaseException: return 3
    except ZeroDivisionError: return 8
    else: return 34
    return res

print(f())
\end{verbatim}


\vspace{5mm}
               
{\begin{center}Завдання {3}\end{center}}
\vspace{-6mm}

\begin{verbatim}
j=7
while j>6:
    j+=-3
print(j)
\end{verbatim}


\end{minipage}
\vline \hspace{6mm}
\begin{minipage}[t]{10.7cm}
               
{\begin{center}Завдання {4}\end{center}}
\vspace{-6mm}

\begin{verbatim}
for a in range(5,5+3,3):
    if a<=4:
        continue
    print(a, end=' ')
    a=0
else:
    print(a, end=' ')
print(a, end=' ')
\end{verbatim}


\vspace{5mm}
               
{\begin{center}Завдання {5}\end{center}}
\vspace{-6mm}

\begin{verbatim}
def f(n):
    print(n % 10, end=' ')
    if n > 9:
        f(n // 100)
    

f(987654)
\end{verbatim}


\end{minipage}

\newpage

{{22.11.2024 Контрольна робота \No 2, варіант  \No \textbf { 43 }\par}}
{
%\begin {large}

Тестова частина. Завдання 1-5: Що буде виведено за виконання (повідомлення інтерпретатора про помилку позначати error)?

Задачі (на звороті): написати функцію для розв'язання задачі.

\vspace{5mm} \begin{minipage}[t]{8.5cm}
               
{\begin{center}Завдання {1}\end{center}}
\vspace{-6mm}

\begin{verbatim}
a,b,c=4,0,6
def f(a):
    global c
    a=4
    b+=3
    c=2
    return a+b+c

a,b,c=2,7,5
print(f(a),a,b,c)
\end{verbatim}


\vspace{5mm}
               
{\begin{center}Завдання {2}\end{center}}
\vspace{-6mm}

\begin{verbatim}
def f():
    try:
        res = 6<4
    except KeyboardInterrupt: return 2
    except Exception: return 0
    return res

print(f())
\end{verbatim}


\vspace{5mm}
               
{\begin{center}Завдання {3}\end{center}}
\vspace{-6mm}

\begin{verbatim}
t=6
while t<=9:
    t+=2
print(t)
\end{verbatim}


\end{minipage}
\vline \hspace{6mm}
\begin{minipage}[t]{10.7cm}
               
{\begin{center}Завдання {4}\end{center}}
\vspace{-6mm}

\begin{verbatim}
for a in range(-2,-2+6,-2):
    if a<3:
        break
    print(a, end=' ')
    a=-10
else:
    print(13, end=' ')
print(a, end=' ')
\end{verbatim}


\vspace{5mm}
               
{\begin{center}Завдання {5}\end{center}}
\vspace{-6mm}

\begin{verbatim}
def f(n):
    if n > 9:
        return f(n // 100) + 1
    else:
        return 1

print(f(123456))
\end{verbatim}


\end{minipage}

\newpage

{{22.11.2024 Контрольна робота \No 2, варіант  \No \textbf { 44 }\par}}
{
%\begin {large}

Тестова частина. Завдання 1-5: Що буде виведено за виконання (повідомлення інтерпретатора про помилку позначати error)?

Задачі (на звороті): написати функцію для розв'язання задачі.

\vspace{5mm} \begin{minipage}[t]{8.5cm}
               
{\begin{center}Завдання {1}\end{center}}
\vspace{-6mm}

\begin{verbatim}
a,b,c=2,9,4
def h(b):
    a-=3
    b=4
    c=5
    return a+b+c

a,b,c=6,3,8
print(h(a),a,b,c)
\end{verbatim}


\vspace{5mm}
               
{\begin{center}Завдання {2}\end{center}}
\vspace{-6mm}

\begin{verbatim}
def f():
    try:
        res = int("1")
        return 41
    except TypeError: return 2
    except ValueError: return 7
    else: return 31
    finally: return 21
    return res

print(f())
\end{verbatim}


\vspace{5mm}
               
{\begin{center}Завдання {3}\end{center}}
\vspace{-6mm}

\begin{verbatim}
n=14
while n>=1:
    n+=-3
print(n)
\end{verbatim}


\end{minipage}
\vline \hspace{6mm}
\begin{minipage}[t]{10.7cm}
               
{\begin{center}Завдання {4}\end{center}}
\vspace{-6mm}

\begin{verbatim}
for b in range(-7,-7+5,-3):
    if b>4:
        continue
    print(b, end=' ')
    b=2
    if b<7:
        break
else:
    print(b, end=' ')
print(b, end=' ')
\end{verbatim}


\vspace{5mm}
               
{\begin{center}Завдання {5}\end{center}}
\vspace{-6mm}

\begin{verbatim}
def f(n):
    if n > 9:
        f(n // 10)
    print(n % 10, end=' ')
    

f(123456)
\end{verbatim}


\end{minipage}

\newpage

{{22.11.2024 Контрольна робота \No 2, варіант  \No \textbf { 45 }\par}}
{
%\begin {large}

Тестова частина. Завдання 1-5: Що буде виведено за виконання (повідомлення інтерпретатора про помилку позначати error)?

Задачі (на звороті): написати функцію для розв'язання задачі.

\vspace{5mm} \begin{minipage}[t]{8.5cm}
               
{\begin{center}Завдання {1}\end{center}}
\vspace{-6mm}

\begin{verbatim}
a,b,c=4,9,8
def f(a):
    global c
    a=3
    b=5
    c=1
    return a+b+c

a,b,c=7,7,2
print(f(a),a,b,c)
\end{verbatim}


\vspace{5mm}
               
{\begin{center}Завдання {2}\end{center}}
\vspace{-6mm}

\begin{verbatim}
def f():
    try:
        res = int("d4")
        return 43
    except ZeroDivisionError: return 0
    except BaseException: return 9
    else: return 30
    finally: return 25
    return res

print(f())
\end{verbatim}


\vspace{5mm}
               
{\begin{center}Завдання {3}\end{center}}
\vspace{-6mm}

\begin{verbatim}
j=8
while j<5:
    j+=1
print(j)
\end{verbatim}


\end{minipage}
\vline \hspace{6mm}
\begin{minipage}[t]{10.7cm}
               
{\begin{center}Завдання {4}\end{center}}
\vspace{-6mm}

\begin{verbatim}
for f in range(1,1+4,3):
    if f<=8:
        pass
    print(f, end=' ')
    f=4
else:
    print(f, end=' ')
print(f, end=' ')
\end{verbatim}


\vspace{5mm}
               
{\begin{center}Завдання {5}\end{center}}
\vspace{-6mm}

\begin{verbatim}
def f(n):
    print(n % 10, end=' ')
    if n > 9:
        f(n // 10)
    

f(543216)
\end{verbatim}


\end{minipage}

\newpage

{{22.11.2024 Контрольна робота \No 2, варіант  \No \textbf { 46 }\par}}
{
%\begin {large}

Тестова частина. Завдання 1-5: Що буде виведено за виконання (повідомлення інтерпретатора про помилку позначати error)?

Задачі (на звороті): написати функцію для розв'язання задачі.

\vspace{5mm} \begin{minipage}[t]{8.5cm}
               
{\begin{center}Завдання {1}\end{center}}
\vspace{-6mm}

\begin{verbatim}
a,b,c=0,1,7
def g(a):
    global c
    a-=1
    b=4
    c=3
    return a+b+c

a,b,c=6,1,0
print(g(b),a,b,c)
\end{verbatim}


\vspace{5mm}
               
{\begin{center}Завдання {2}\end{center}}
\vspace{-6mm}

\begin{verbatim}
def f():
    try:
        res = int(1%0)
    except TypeError: return 5
    except ZeroDivisionError: return 6
    return res

print(f())
\end{verbatim}


\vspace{5mm}
               
{\begin{center}Завдання {3}\end{center}}
\vspace{-6mm}

\begin{verbatim}
i=6
while i>5:
    i+=-2
print(i)
\end{verbatim}


\end{minipage}
\vline \hspace{6mm}
\begin{minipage}[t]{10.7cm}
               
{\begin{center}Завдання {4}\end{center}}
\vspace{-6mm}

\begin{verbatim}
for d in range(-4,-4+2,-3):
    if d<7:
        continue
    print(d, end=' ')
    d=-1
else:
    print(d, end=' ')
print(d, end=' ')
\end{verbatim}


\vspace{5mm}
               
{\begin{center}Завдання {5}\end{center}}
\vspace{-6mm}

\begin{verbatim}
def f(n):
    if n > 2:
        f(n // 2)
    print(n % 2, end=' ')
    

f(16)
\end{verbatim}


\end{minipage}

\newpage

{{22.11.2024 Контрольна робота \No 2, варіант  \No \textbf { 47 }\par}}
{
%\begin {large}

Тестова частина. Завдання 1-5: Що буде виведено за виконання (повідомлення інтерпретатора про помилку позначати error)?

Задачі (на звороті): написати функцію для розв'язання задачі.

\vspace{5mm} \begin{minipage}[t]{8.5cm}
               
{\begin{center}Завдання {1}\end{center}}
\vspace{-6mm}

\begin{verbatim}
a,b,c=2,8,3
def f(b):
    global c
    a=1
    b+=5
    c=2
    return a+b+c

a,b,c=5,4,9
print(f(b),a,b,c)
\end{verbatim}


\vspace{5mm}
               
{\begin{center}Завдання {2}\end{center}}
\vspace{-6mm}

\begin{verbatim}
def f():
    try:
        res = int(3//2)
        return 44
    except ValueError: return 8
    except BaseException: return 9
    else: return 32
    finally: return 20
    return res

print(f())
\end{verbatim}


\vspace{5mm}
               
{\begin{center}Завдання {3}\end{center}}
\vspace{-6mm}

\begin{verbatim}
m=10
while m>=9:
    m+=-3
print(m)
\end{verbatim}


\end{minipage}
\vline \hspace{6mm}
\begin{minipage}[t]{10.7cm}
               
{\begin{center}Завдання {4}\end{center}}
\vspace{-6mm}

\begin{verbatim}
for e in range(0,0+4,2):
    if e>=6:
        continue
    print(e, end=' ')
    e=5
else:
    print(e, end=' ')
print(e, end=' ')
\end{verbatim}


\vspace{5mm}
               
{\begin{center}Завдання {5}\end{center}}
\vspace{-6mm}

\begin{verbatim}
def f(n):
    if n > 9:
        f(n // 10)
    print(n % 10, end=' ')
    

f(543216)
\end{verbatim}


\end{minipage}

\newpage

{{22.11.2024 Контрольна робота \No 2, варіант  \No \textbf { 48 }\par}}
{
%\begin {large}

Тестова частина. Завдання 1-5: Що буде виведено за виконання (повідомлення інтерпретатора про помилку позначати error)?

Задачі (на звороті): написати функцію для розв'язання задачі.

\vspace{5mm} \begin{minipage}[t]{8.5cm}
               
{\begin{center}Завдання {1}\end{center}}
\vspace{-6mm}

\begin{verbatim}
a,b,c=9,4,6
def h(a):
    a=1
    b*=2
    c=3
    return a+b+c

a,b,c=3,0,7
print(h(a),a,b,c)
\end{verbatim}


\vspace{5mm}
               
{\begin{center}Завдання {2}\end{center}}
\vspace{-6mm}

\begin{verbatim}
def f():
    try:
        res = int("3")
    except Exception: return 2
    except KeyboardInterrupt: return 4
    return res

print(f())
\end{verbatim}


\vspace{5mm}
               
{\begin{center}Завдання {3}\end{center}}
\vspace{-6mm}

\begin{verbatim}
n=7
while n>4:
    n+=-2
print(n)
\end{verbatim}


\end{minipage}
\vline \hspace{6mm}
\begin{minipage}[t]{10.7cm}
               
{\begin{center}Завдання {4}\end{center}}
\vspace{-6mm}

\begin{verbatim}
for c in range(8,8+5,-2):
    if c>5:
        break
    print(c, end=' ')
    c=-4
else:
    print(c, end=' ')
print(c, end=' ')
\end{verbatim}


\vspace{5mm}
               
{\begin{center}Завдання {5}\end{center}}
\vspace{-6mm}

\begin{verbatim}
def f(n):
    if n <= 9:
        print(n % 10, end=' ')
    else:
        f(n // 10)

f(543216)
\end{verbatim}


\end{minipage}

\newpage

{{22.11.2024 Контрольна робота \No 2, варіант  \No \textbf { 49 }\par}}
{
%\begin {large}

Тестова частина. Завдання 1-5: Що буде виведено за виконання (повідомлення інтерпретатора про помилку позначати error)?

Задачі (на звороті): написати функцію для розв'язання задачі.

\vspace{5mm} \begin{minipage}[t]{8.5cm}
               
{\begin{center}Завдання {1}\end{center}}
\vspace{-6mm}

\begin{verbatim}
a,b,c=2,8,5
def g(b):
    global c
    a-=5
    b=4
    c=5
    return a+b+c

a,b,c=1,8,7
print(g(b),a,b,c)
\end{verbatim}


\vspace{5mm}
               
{\begin{center}Завдання {2}\end{center}}
\vspace{-6mm}

\begin{verbatim}
def f():
    try:
        res = 7>1
    except ZeroDivisionError: return 8
    except ValueError: return 0
    else: return 31
    finally: return 20
    return res

print(f())
\end{verbatim}


\vspace{5mm}
               
{\begin{center}Завдання {3}\end{center}}
\vspace{-6mm}

\begin{verbatim}
i=5
while i>=3:
    i+=-3
print(i)
\end{verbatim}


\end{minipage}
\vline \hspace{6mm}
\begin{minipage}[t]{10.7cm}
               
{\begin{center}Завдання {4}\end{center}}
\vspace{-6mm}

\begin{verbatim}
for a in range(-3,-3+3,-2):
    if a<=3:
        break
    print(a, end=' ')
    a=-7
else:
    print(a, end=' ')
print(a, end=' ')
\end{verbatim}


\vspace{5mm}
               
{\begin{center}Завдання {5}\end{center}}
\vspace{-6mm}

\begin{verbatim}
def f(n):
    if n > 9:
        f(n // 100)
    else:
        print(n % 10, end=' ')

f(123456)
\end{verbatim}


\end{minipage}

\newpage

{{22.11.2024 Контрольна робота \No 2, варіант  \No \textbf { 50 }\par}}
{
%\begin {large}

Тестова частина. Завдання 1-5: Що буде виведено за виконання (повідомлення інтерпретатора про помилку позначати error)?

Задачі (на звороті): написати функцію для розв'язання задачі.

\vspace{5mm} \begin{minipage}[t]{8.5cm}
               
{\begin{center}Завдання {1}\end{center}}
\vspace{-6mm}

\begin{verbatim}
a,b,c=0,4,1
def h(b):
    a=4
    b=2
    c=3
    return a+b+c

a,b,c=8,9,6
print(h(a),a,b,c)
\end{verbatim}


\vspace{5mm}
               
{\begin{center}Завдання {2}\end{center}}
\vspace{-6mm}

\begin{verbatim}
def f():
    try:
        res = int("b5")
        return 43
    except Exception: return 1
    except BaseException: return 6
    return res

print(f())
\end{verbatim}


\vspace{5mm}
               
{\begin{center}Завдання {3}\end{center}}
\vspace{-6mm}

\begin{verbatim}
t=15
while t>8:
    t+=-2
print(t)
\end{verbatim}


\end{minipage}
\vline \hspace{6mm}
\begin{minipage}[t]{10.7cm}
               
{\begin{center}Завдання {4}\end{center}}
\vspace{-6mm}

\begin{verbatim}
for c in range(5,5+6,2):
    if c>8:
        continue
    print(c, end=' ')
    c=7
else:
    print(c, end=' ')
print(c, end=' ')
\end{verbatim}


\vspace{5mm}
               
{\begin{center}Завдання {5}\end{center}}
\vspace{-6mm}

\begin{verbatim}
def f(n):
    if n > 2:
        f(n // 2)
    else:
        print(n % 2, end=' ')

f(16)
\end{verbatim}


\end{minipage}

\newpage

{{22.11.2024 Контрольна робота \No 2, варіант  \No \textbf { 51 }\par}}
{
%\begin {large}

Тестова частина. Завдання 1-5: Що буде виведено за виконання (повідомлення інтерпретатора про помилку позначати error)?

Задачі (на звороті): написати функцію для розв'язання задачі.

\vspace{5mm} \begin{minipage}[t]{8.5cm}
               
{\begin{center}Завдання {1}\end{center}}
\vspace{-6mm}

\begin{verbatim}
a,b,c=9,0,3
def h(a):
    a=4
    b*=1
    c=3
    return a+b+c

a,b,c=2,5,1
print(h(b),a,b,c)
\end{verbatim}


\vspace{5mm}
               
{\begin{center}Завдання {2}\end{center}}
\vspace{-6mm}

\begin{verbatim}
def f():
    try:
        res = int(1/2)
        return 40
    except TypeError: return 6
    except KeyboardInterrupt: return 9
    else: return 35
    return res

print(f())
\end{verbatim}


\vspace{5mm}
               
{\begin{center}Завдання {3}\end{center}}
\vspace{-6mm}

\begin{verbatim}
k=11
while k>=7:
    k+=-2
print(k)
\end{verbatim}


\end{minipage}
\vline \hspace{6mm}
\begin{minipage}[t]{10.7cm}
               
{\begin{center}Завдання {4}\end{center}}
\vspace{-6mm}

\begin{verbatim}
for e in range(6,6+5,3):
    if e<4:
        continue
    print(e, end=' ')
    e=8
else:
    print(e, end=' ')
print(e, end=' ')
\end{verbatim}


\vspace{5mm}
               
{\begin{center}Завдання {5}\end{center}}
\vspace{-6mm}

\begin{verbatim}
def f(n):
    if n > 9:
        return f(n // 10) + 1
    else:
        return 1

print(f(123456))
\end{verbatim}


\end{minipage}

\newpage

{{22.11.2024 Контрольна робота \No 2, варіант  \No \textbf { 52 }\par}}
{
%\begin {large}

Тестова частина. Завдання 1-5: Що буде виведено за виконання (повідомлення інтерпретатора про помилку позначати error)?

Задачі (на звороті): написати функцію для розв'язання задачі.

\vspace{5mm} \begin{minipage}[t]{8.5cm}
               
{\begin{center}Завдання {1}\end{center}}
\vspace{-6mm}

\begin{verbatim}
a,b,c=6,4,3
def f(a):
    global c
    a=2
    b+=2
    c=1
    return a+b+c

a,b,c=8,1,6
print(f(b),a,b,c)
\end{verbatim}


\vspace{5mm}
               
{\begin{center}Завдання {2}\end{center}}
\vspace{-6mm}

\begin{verbatim}
def f():
    try:
        res = int("c2")
    except Exception: return 0
    except ZeroDivisionError: return 3
    finally: return 21
    return res

print(f())
\end{verbatim}


\vspace{5mm}
               
{\begin{center}Завдання {3}\end{center}}
\vspace{-6mm}

\begin{verbatim}
n=4
while n<12:
    n+=1
print(n)
\end{verbatim}


\end{minipage}
\vline \hspace{6mm}
\begin{minipage}[t]{10.7cm}
               
{\begin{center}Завдання {4}\end{center}}
\vspace{-6mm}

\begin{verbatim}
for f in range(4,4+4,-3):
    if f>=5:
        pass
    print(f, end=' ')
    f=10
else:
    print(13, end=' ')
print(f, end=' ')
\end{verbatim}


\vspace{5mm}
               
{\begin{center}Завдання {5}\end{center}}
\vspace{-6mm}

\begin{verbatim}
def f(n):
    print(n % 10, end=' ')
    if n > 9:
        f(n // 10)
    

f(123456)
\end{verbatim}


\end{minipage}

\newpage

{{22.11.2024 Контрольна робота \No 2, варіант  \No \textbf { 53 }\par}}
{
%\begin {large}

Тестова частина. Завдання 1-5: Що буде виведено за виконання (повідомлення інтерпретатора про помилку позначати error)?

Задачі (на звороті): написати функцію для розв'язання задачі.

\vspace{5mm} \begin{minipage}[t]{8.5cm}
               
{\begin{center}Завдання {1}\end{center}}
\vspace{-6mm}

\begin{verbatim}
a,b,c=4,7,5
def h(b):
    a+=3
    b=4
    c=5
    return a+b+c

a,b,c=0,2,9
print(h(a),a,b,c)
\end{verbatim}


\vspace{5mm}
               
{\begin{center}Завдання {2}\end{center}}
\vspace{-6mm}

\begin{verbatim}
def f():
    try:
        res = 5!=2
    except BaseException: return 4
    except KeyboardInterrupt: return 7
    finally: return 25
    return res

print(f())
\end{verbatim}


\vspace{5mm}
               
{\begin{center}Завдання {3}\end{center}}
\vspace{-6mm}

\begin{verbatim}
m=7
while m<=11:
    m+=2
print(m)
\end{verbatim}


\end{minipage}
\vline \hspace{6mm}
\begin{minipage}[t]{10.7cm}
               
{\begin{center}Завдання {4}\end{center}}
\vspace{-6mm}

\begin{verbatim}
for d in range(-6,-6+3,-2):
    if d<=7:
        continue
    print(d, end=' ')
    d=0
else:
    print(d, end=' ')
print(d, end=' ')
\end{verbatim}


\vspace{5mm}
               
{\begin{center}Завдання {5}\end{center}}
\vspace{-6mm}

\begin{verbatim}
def f(n):
    if n <= 9:
        print(n % 10, end=' ')
    else:
        f(n // 10)

f(987651)
\end{verbatim}


\end{minipage}

\newpage

{{22.11.2024 Контрольна робота \No 2, варіант  \No \textbf { 54 }\par}}
{
%\begin {large}

Тестова частина. Завдання 1-5: Що буде виведено за виконання (повідомлення інтерпретатора про помилку позначати error)?

Задачі (на звороті): написати функцію для розв'язання задачі.

\vspace{5mm} \begin{minipage}[t]{8.5cm}
               
{\begin{center}Завдання {1}\end{center}}
\vspace{-6mm}

\begin{verbatim}
a,b,c=3,5,2
def f(a):
    global c
    a=5
    b=5
    c=4
    return a+b+c

a,b,c=1,7,3
print(f(a),a,b,c)
\end{verbatim}


\vspace{5mm}
               
{\begin{center}Завдання {2}\end{center}}
\vspace{-6mm}

\begin{verbatim}
def f():
    try:
        res = int("8")
        return 41
    except TypeError: return 9
    except ValueError: return 7
    else: return 34
    return res

print(f())
\end{verbatim}


\vspace{5mm}
               
{\begin{center}Завдання {3}\end{center}}
\vspace{-6mm}

\begin{verbatim}
t=2
while t<=10:
    t+=2
print(t)
\end{verbatim}


\end{minipage}
\vline \hspace{6mm}
\begin{minipage}[t]{10.7cm}
               
{\begin{center}Завдання {4}\end{center}}
\vspace{-6mm}

\begin{verbatim}
for b in range(-8,-8+2,-3):
    if b>6:
        continue
    print(b, end=' ')
    b=9
else:
    print(b, end=' ')
print(b, end=' ')
\end{verbatim}


\vspace{5mm}
               
{\begin{center}Завдання {5}\end{center}}
\vspace{-6mm}

\begin{verbatim}
def f(n):
    if n <= 9:
        print(n % 10, end=' ')
    else:
        f(n // 10)

f(123456)
\end{verbatim}


\end{minipage}

\newpage

{{22.11.2024 Контрольна робота \No 2, варіант  \No \textbf { 55 }\par}}
{
%\begin {large}

Тестова частина. Завдання 1-5: Що буде виведено за виконання (повідомлення інтерпретатора про помилку позначати error)?

Задачі (на звороті): написати функцію для розв'язання задачі.

\vspace{5mm} \begin{minipage}[t]{8.5cm}
               
{\begin{center}Завдання {1}\end{center}}
\vspace{-6mm}

\begin{verbatim}
a,b,c=0,8,6
def g(b):
    global c
    a=3
    b=2
    c=1
    return a+b+c

a,b,c=9,4,5
print(g(b),a,b,c)
\end{verbatim}


\vspace{5mm}
               
{\begin{center}Завдання {2}\end{center}}
\vspace{-6mm}

\begin{verbatim}
def f():
    try:
        res = int(2/0.0)
    except TypeError: return 0
    except BaseException: return 1
    else: return 32
    return res

print(f())
\end{verbatim}


\vspace{5mm}
               
{\begin{center}Завдання {3}\end{center}}
\vspace{-6mm}

\begin{verbatim}
i=0
while i<8:
    i+=1
print(i)
\end{verbatim}


\end{minipage}
\vline \hspace{6mm}
\begin{minipage}[t]{10.7cm}
               
{\begin{center}Завдання {4}\end{center}}
\vspace{-6mm}

\begin{verbatim}
for b in range(10,10+6,3):
    if b>=6:
        break
    print(b, end=' ')
    b=-5
    if b>=4:
        break
else:
    print(b, end=' ')
print(b, end=' ')
\end{verbatim}


\vspace{5mm}
               
{\begin{center}Завдання {5}\end{center}}
\vspace{-6mm}

\begin{verbatim}
def f(n):
    if n > 9:
            f(n // 100)
    print(n % 10, end=' ')
    

f(987654)
\end{verbatim}


\end{minipage}

\newpage

{{22.11.2024 Контрольна робота \No 2, варіант  \No \textbf { 56 }\par}}
{
%\begin {large}

Тестова частина. Завдання 1-5: Що буде виведено за виконання (повідомлення інтерпретатора про помилку позначати error)?

Задачі (на звороті): написати функцію для розв'язання задачі.

\vspace{5mm} \begin{minipage}[t]{8.5cm}
               
{\begin{center}Завдання {1}\end{center}}
\vspace{-6mm}

\begin{verbatim}
a,b,c=2,7,8
def g(a):
    a*=3
    b=1
    c=2
    return a+b+c

a,b,c=1,9,5
print(g(a),a,b,c)
\end{verbatim}


\vspace{5mm}
               
{\begin{center}Завдання {2}\end{center}}
\vspace{-6mm}

\begin{verbatim}
def f():
    try:
        res = int(4%1)
        return 45
    except ZeroDivisionError: return 8
    except Exception: return 3
    finally: return 22
    return res

print(f())
\end{verbatim}


\vspace{5mm}
               
{\begin{center}Завдання {3}\end{center}}
\vspace{-6mm}

\begin{verbatim}
m=9
while m<=7:
    m+=2
print(m)
\end{verbatim}


\end{minipage}
\vline \hspace{6mm}
\begin{minipage}[t]{10.7cm}
               
{\begin{center}Завдання {4}\end{center}}
\vspace{-6mm}

\begin{verbatim}
for d in range(-9,-9+5,2):
    if d<7:
        pass
    print(d, end=' ')
    d=-3
else:
    print(13, end=' ')
print(d, end=' ')
\end{verbatim}


\vspace{5mm}
               
{\begin{center}Завдання {5}\end{center}}
\vspace{-6mm}

\begin{verbatim}
def f(n):
    if n > 9:
        f(n // 100)
    print(n % 10, end=' ')
    

f(123456)
\end{verbatim}


\end{minipage}

\newpage

{{22.11.2024 Контрольна робота \No 2, варіант  \No \textbf { 57 }\par}}
{
%\begin {large}

Тестова частина. Завдання 1-5: Що буде виведено за виконання (повідомлення інтерпретатора про помилку позначати error)?

Задачі (на звороті): написати функцію для розв'язання задачі.

\vspace{5mm} \begin{minipage}[t]{8.5cm}
               
{\begin{center}Завдання {1}\end{center}}
\vspace{-6mm}

\begin{verbatim}
a,b,c=0,6,3
def f(b):
    global c
    a-=4
    b=5
    c=4
    return a+b+c

a,b,c=4,2,0
print(f(a),a,b,c)
\end{verbatim}


\vspace{5mm}
               
{\begin{center}Завдання {2}\end{center}}
\vspace{-6mm}

\begin{verbatim}
def f():
    try:
        res = int("a5")
        return 42
    except ValueError: return 6
    except KeyboardInterrupt: return 3
    finally: return 23
    return res

print(f())
\end{verbatim}


\vspace{5mm}
               
{\begin{center}Завдання {3}\end{center}}
\vspace{-6mm}

\begin{verbatim}
j=3
while j<7:
    j+=1
print(j)
\end{verbatim}


\end{minipage}
\vline \hspace{6mm}
\begin{minipage}[t]{10.7cm}
               
{\begin{center}Завдання {4}\end{center}}
\vspace{-6mm}

\begin{verbatim}
for a in range(-5,-5+2,2):
    if a<5:
        continue
    print(a, end=' ')
    a=6
else:
    print(a, end=' ')
print(a, end=' ')
\end{verbatim}


\vspace{5mm}
               
{\begin{center}Завдання {5}\end{center}}
\vspace{-6mm}

\begin{verbatim}
def f(n):
    if n > 9:
        f(n // 10)
    else:
        print(n % 10, end=' ')

f(123456)
\end{verbatim}


\end{minipage}

\newpage

{{22.11.2024 Контрольна робота \No 2, варіант  \No \textbf { 58 }\par}}
{
%\begin {large}

Тестова частина. Завдання 1-5: Що буде виведено за виконання (повідомлення інтерпретатора про помилку позначати error)?

Задачі (на звороті): написати функцію для розв'язання задачі.

\vspace{5mm} \begin{minipage}[t]{8.5cm}
               
{\begin{center}Завдання {1}\end{center}}
\vspace{-6mm}

\begin{verbatim}
a,b,c=1,9,5
def f(b):
    global c
    a+=3
    b=1
    c=5
    return a+b+c

a,b,c=3,8,4
print(f(b),a,b,c)
\end{verbatim}


\vspace{5mm}
               
{\begin{center}Завдання {2}\end{center}}
\vspace{-6mm}

\begin{verbatim}
def f():
    try:
        res = int(7//0)
    except KeyboardInterrupt: return 5
    except ValueError: return 6
    else: return 30
    return res

print(f())
\end{verbatim}


\vspace{5mm}
               
{\begin{center}Завдання {3}\end{center}}
\vspace{-6mm}

\begin{verbatim}
k=5
while k<11:
    k+=2
print(k)
\end{verbatim}


\end{minipage}
\vline \hspace{6mm}
\begin{minipage}[t]{10.7cm}
               
{\begin{center}Завдання {4}\end{center}}
\vspace{-6mm}

\begin{verbatim}
for e in range(7,7+3,-2):
    if e>=4:
        pass
    print(e, end=' ')
    e=-2
    if e>8:
        break
else:
    print(e, end=' ')
print(e, end=' ')
\end{verbatim}


\vspace{5mm}
               
{\begin{center}Завдання {5}\end{center}}
\vspace{-6mm}

\begin{verbatim}
def f(n):
    if n > 9:
        f(n // 10)
    else:
        print(n % 10, end=' ')

f(543216)
\end{verbatim}


\end{minipage}

\newpage

{{22.11.2024 Контрольна робота \No 2, варіант  \No \textbf { 59 }\par}}
{
%\begin {large}

Тестова частина. Завдання 1-5: Що буде виведено за виконання (повідомлення інтерпретатора про помилку позначати error)?

Задачі (на звороті): написати функцію для розв'язання задачі.

\vspace{5mm} \begin{minipage}[t]{8.5cm}
               
{\begin{center}Завдання {1}\end{center}}
\vspace{-6mm}

\begin{verbatim}
a,b,c=6,7,7
def h(a):
    a=2
    b*=3
    c=4
    return a+b+c

a,b,c=8,9,0
print(h(b),a,b,c)
\end{verbatim}


\vspace{5mm}
               
{\begin{center}Завдання {2}\end{center}}
\vspace{-6mm}

\begin{verbatim}
def f():
    try:
        res = int("2")
        return 44
    except ZeroDivisionError: return 9
    except BaseException: return 8
    return res

print(f())
\end{verbatim}


\vspace{5mm}
               
{\begin{center}Завдання {3}\end{center}}
\vspace{-6mm}

\begin{verbatim}
n=1
while n<=6:
    n+=1
print(n)
\end{verbatim}


\end{minipage}
\vline \hspace{6mm}
\begin{minipage}[t]{10.7cm}
               
{\begin{center}Завдання {4}\end{center}}
\vspace{-6mm}

\begin{verbatim}
for c in range(-1,-1+6,-3):
    if c>8:
        break
    print(c, end=' ')
    c=-8
else:
    print(c, end=' ')
print(c, end=' ')
\end{verbatim}


\vspace{5mm}
               
{\begin{center}Завдання {5}\end{center}}
\vspace{-6mm}

\begin{verbatim}
def f(n):
    print(n % 2, end=' ')
    if n > 2:
        f(n // 2)
    

f(16)
\end{verbatim}


\end{minipage}

\newpage

{{22.11.2024 Контрольна робота \No 2, варіант  \No \textbf { 60 }\par}}
{
%\begin {large}

Тестова частина. Завдання 1-5: Що буде виведено за виконання (повідомлення інтерпретатора про помилку позначати error)?

Задачі (на звороті): написати функцію для розв'язання задачі.

\vspace{5mm} \begin{minipage}[t]{8.5cm}
               
{\begin{center}Завдання {1}\end{center}}
\vspace{-6mm}

\begin{verbatim}
a,b,c=2,4,6
def g(b):
    global c
    a=5
    b-=2
    c=1
    return a+b+c

a,b,c=5,3,1
print(g(a),a,b,c)
\end{verbatim}


\vspace{5mm}
               
{\begin{center}Завдання {2}\end{center}}
\vspace{-6mm}

\begin{verbatim}
def f():
    try:
        res = 4==9
    except TypeError: return 0
    except Exception: return 1
    else: return 33
    finally: return 24
    return res

print(f())
\end{verbatim}


\vspace{5mm}
               
{\begin{center}Завдання {3}\end{center}}
\vspace{-6mm}

\begin{verbatim}
j=8
while j>0:
    j+=-3
print(j)
\end{verbatim}


\end{minipage}
\vline \hspace{6mm}
\begin{minipage}[t]{10.7cm}
               
{\begin{center}Завдання {4}\end{center}}
\vspace{-6mm}

\begin{verbatim}
for f in range(2,2+4,3):
    if f<=3:
        continue
    print(f, end=' ')
    f=-9
else:
    print(f, end=' ')
print(f, end=' ')
\end{verbatim}


\vspace{5mm}
               
{\begin{center}Завдання {5}\end{center}}
\vspace{-6mm}

\begin{verbatim}
def f(n):
    if n > 9:
        f(n // 100)
    else:
        print(n % 10, end=' ')

f(987654)
\end{verbatim}


\end{minipage}

\newpage

{{22.11.2024 Контрольна робота \No 2, варіант  \No \textbf { 61 }\par}}
{
%\begin {large}

Тестова частина. Завдання 1-5: Що буде виведено за виконання (повідомлення інтерпретатора про помилку позначати error)?

Задачі (на звороті): написати функцію для розв'язання задачі.

\vspace{5mm} \begin{minipage}[t]{8.5cm}
               
{\begin{center}Завдання {1}\end{center}}
\vspace{-6mm}

\begin{verbatim}
a,b,c=5,2,4
def f(a):
    a+=5
    b=4
    c=1
    return a+b+c

a,b,c=9,1,0
print(f(a),a,b,c)
\end{verbatim}


\vspace{5mm}
               
{\begin{center}Завдання {2}\end{center}}
\vspace{-6mm}

\begin{verbatim}
def f():
    try:
        res = 0>8
    except TypeError: return 2
    except ZeroDivisionError: return 4
    else: return 34
    return res

print(f())
\end{verbatim}


\vspace{5mm}
               
{\begin{center}Завдання {3}\end{center}}
\vspace{-6mm}

\begin{verbatim}
i=9
while i<=13:
    i+=1
print(i)
\end{verbatim}


\end{minipage}
\vline \hspace{6mm}
\begin{minipage}[t]{10.7cm}
               
{\begin{center}Завдання {4}\end{center}}
\vspace{-6mm}

\begin{verbatim}
for c in range(-10,-10+3,-3):
    if c>=4:
        pass
    print(c, end=' ')
    c=-6
else:
    print(c, end=' ')
print(c, end=' ')
\end{verbatim}


\vspace{5mm}
               
{\begin{center}Завдання {5}\end{center}}
\vspace{-6mm}

\begin{verbatim}
def f(n):
    if n > 2:
        f(n // 2)
    print(n % 2, end=' ')
    

f(16)
\end{verbatim}


\end{minipage}

\newpage

{{22.11.2024 Контрольна робота \No 2, варіант  \No \textbf { 62 }\par}}
{
%\begin {large}

Тестова частина. Завдання 1-5: Що буде виведено за виконання (повідомлення інтерпретатора про помилку позначати error)?

Задачі (на звороті): написати функцію для розв'язання задачі.

\vspace{5mm} \begin{minipage}[t]{8.5cm}
               
{\begin{center}Завдання {1}\end{center}}
\vspace{-6mm}

\begin{verbatim}
a,b,c=6,8,7
def h(b):
    global c
    a=2
    b-=3
    c=5
    return a+b+c

a,b,c=3,7,4
print(h(a),a,b,c)
\end{verbatim}


\vspace{5mm}
               
{\begin{center}Завдання {2}\end{center}}
\vspace{-6mm}

\begin{verbatim}
def f():
    try:
        res = int("c3")
        return 43
    except Exception: return 1
    except ValueError: return 6
    finally: return 20
    return res

print(f())
\end{verbatim}


\vspace{5mm}
               
{\begin{center}Завдання {3}\end{center}}
\vspace{-6mm}

\begin{verbatim}
m=6
while m<8:
    m+=2
print(m)
\end{verbatim}


\end{minipage}
\vline \hspace{6mm}
\begin{minipage}[t]{10.7cm}
               
{\begin{center}Завдання {4}\end{center}}
\vspace{-6mm}

\begin{verbatim}
for d in range(3,3+5,-2):
    if d>5:
        break
    print(d, end=' ')
    d=1
else:
    print(d, end=' ')
print(d, end=' ')
\end{verbatim}


\vspace{5mm}
               
{\begin{center}Завдання {5}\end{center}}
\vspace{-6mm}

\begin{verbatim}
def f(n):
    if n > 9:
        f(n // 100)
    else:
        print(n % 10, end=' ')

f(123456)
\end{verbatim}


\end{minipage}

\newpage

{{22.11.2024 Контрольна робота \No 2, варіант  \No \textbf { 63 }\par}}
{
%\begin {large}

Тестова частина. Завдання 1-5: Що буде виведено за виконання (повідомлення інтерпретатора про помилку позначати error)?

Задачі (на звороті): написати функцію для розв'язання задачі.

\vspace{5mm} \begin{minipage}[t]{8.5cm}
               
{\begin{center}Завдання {1}\end{center}}
\vspace{-6mm}

\begin{verbatim}
a,b,c=1,2,9
def g(a):
    a=1
    b*=2
    c=3
    return a+b+c

a,b,c=6,0,3
print(g(b),a,b,c)
\end{verbatim}


\vspace{5mm}
               
{\begin{center}Завдання {2}\end{center}}
\vspace{-6mm}

\begin{verbatim}
def f():
    try:
        res = int(8/0.0)
    except BaseException: return 7
    except KeyboardInterrupt: return 5
    finally: return 25
    return res

print(f())
\end{verbatim}


\vspace{5mm}
               
{\begin{center}Завдання {3}\end{center}}
\vspace{-6mm}

\begin{verbatim}
t=12
while t>=7:
    t+=-2
print(t)
\end{verbatim}


\end{minipage}
\vline \hspace{6mm}
\begin{minipage}[t]{10.7cm}
               
{\begin{center}Завдання {4}\end{center}}
\vspace{-6mm}

\begin{verbatim}
for b in range(9,9+6,2):
    if b<7:
        continue
    print(b, end=' ')
    b=3
else:
    print(b, end=' ')
print(b, end=' ')
\end{verbatim}


\vspace{5mm}
               
{\begin{center}Завдання {5}\end{center}}
\vspace{-6mm}

\begin{verbatim}
def f(n):
    if n > 9:
        f(n // 10)
    print(n % 10, end=' ')
    

f(123456)
\end{verbatim}


\end{minipage}

\newpage

{{22.11.2024 Контрольна робота \No 2, варіант  \No \textbf { 64 }\par}}
{
%\begin {large}

Тестова частина. Завдання 1-5: Що буде виведено за виконання (повідомлення інтерпретатора про помилку позначати error)?

Задачі (на звороті): написати функцію для розв'язання задачі.

\vspace{5mm} \begin{minipage}[t]{8.5cm}
               
{\begin{center}Завдання {1}\end{center}}
\vspace{-6mm}

\begin{verbatim}
a,b,c=8,5,3
def h(b):
    global c
    a-=4
    b=2
    c=3
    return a+b+c

a,b,c=6,7,1
print(h(a),a,b,c)
\end{verbatim}


\vspace{5mm}
               
{\begin{center}Завдання {2}\end{center}}
\vspace{-6mm}

\begin{verbatim}
def f():
    try:
        res = int(9//2)
        return 45
    except ZeroDivisionError: return 9
    except BaseException: return 0
    else: return 33
    return res

print(f())
\end{verbatim}


\vspace{5mm}
               
{\begin{center}Завдання {3}\end{center}}
\vspace{-6mm}

\begin{verbatim}
j=2
while j<12:
    j+=2
print(j)
\end{verbatim}


\end{minipage}
\vline \hspace{6mm}
\begin{minipage}[t]{10.7cm}
               
{\begin{center}Завдання {4}\end{center}}
\vspace{-6mm}

\begin{verbatim}
for a in range(5,5+4,3):
    if a<=8:
        continue
    print(a, end=' ')
    a=-1
else:
    print(13, end=' ')
print(a, end=' ')
\end{verbatim}


\vspace{5mm}
               
{\begin{center}Завдання {5}\end{center}}
\vspace{-6mm}

\begin{verbatim}
def f(n):
    if n > 9:
        return f(n // 100) + 1
    else:
        return 1

print(f(123456))
\end{verbatim}


\end{minipage}

\newpage

{{22.11.2024 Контрольна робота \No 2, варіант  \No \textbf { 65 }\par}}
{
%\begin {large}

Тестова частина. Завдання 1-5: Що буде виведено за виконання (повідомлення інтерпретатора про помилку позначати error)?

Задачі (на звороті): написати функцію для розв'язання задачі.

\vspace{5mm} \begin{minipage}[t]{8.5cm}
               
{\begin{center}Завдання {1}\end{center}}
\vspace{-6mm}

\begin{verbatim}
a,b,c=5,4,3
def h(a):
    a=1
    b=2
    c=5
    return a+b+c

a,b,c=7,0,9
print(h(a),a,b,c)
\end{verbatim}


\vspace{5mm}
               
{\begin{center}Завдання {2}\end{center}}
\vspace{-6mm}

\begin{verbatim}
def f():
    try:
        res = int("6")
        return 44
    except TypeError: return 3
    except Exception: return 4
    else: return 35
    return res

print(f())
\end{verbatim}


\vspace{5mm}
               
{\begin{center}Завдання {3}\end{center}}
\vspace{-6mm}

\begin{verbatim}
i=13
while i>2:
    i+=-3
print(i)
\end{verbatim}


\end{minipage}
\vline \hspace{6mm}
\begin{minipage}[t]{10.7cm}
               
{\begin{center}Завдання {4}\end{center}}
\vspace{-6mm}

\begin{verbatim}
for f in range(-2,-2+2,-3):
    if f>3:
        continue
    print(f, end=' ')
    f=-9
else:
    print(f, end=' ')
print(f, end=' ')
\end{verbatim}


\vspace{5mm}
               
{\begin{center}Завдання {5}\end{center}}
\vspace{-6mm}

\begin{verbatim}
def f(n):
    if n <= 9:
        print(n % 10, end=' ')
    else:
        f(n // 10)

f(543216)
\end{verbatim}


\end{minipage}

\newpage

{{22.11.2024 Контрольна робота \No 2, варіант  \No \textbf { 66 }\par}}
{
%\begin {large}

Тестова частина. Завдання 1-5: Що буде виведено за виконання (повідомлення інтерпретатора про помилку позначати error)?

Задачі (на звороті): написати функцію для розв'язання задачі.

\vspace{5mm} \begin{minipage}[t]{8.5cm}
               
{\begin{center}Завдання {1}\end{center}}
\vspace{-6mm}

\begin{verbatim}
a,b,c=9,2,5
def f(a):
    a=4
    b*=1
    c=5
    return a+b+c

a,b,c=0,4,8
print(f(a),a,b,c)
\end{verbatim}


\vspace{5mm}
               
{\begin{center}Завдання {2}\end{center}}
\vspace{-6mm}

\begin{verbatim}
def f():
    try:
        res = int(8%1)
    except ValueError: return 1
    except KeyboardInterrupt: return 5
    finally: return 23
    return res

print(f())
\end{verbatim}


\vspace{5mm}
               
{\begin{center}Завдання {3}\end{center}}
\vspace{-6mm}

\begin{verbatim}
k=7
while k<=10:
    k+=1
print(k)
\end{verbatim}


\end{minipage}
\vline \hspace{6mm}
\begin{minipage}[t]{10.7cm}
               
{\begin{center}Завдання {4}\end{center}}
\vspace{-6mm}

\begin{verbatim}
for e in range(1,1+2,3):
    if e<=6:
        continue
    print(e, end=' ')
    e=-3
else:
    print(e, end=' ')
print(e, end=' ')
\end{verbatim}


\vspace{5mm}
               
{\begin{center}Завдання {5}\end{center}}
\vspace{-6mm}

\begin{verbatim}
def f(n):
    if n > 9:
        f(n // 10)
    print(n % 10, end=' ')
    

f(543216)
\end{verbatim}


\end{minipage}

\newpage

{{22.11.2024 Контрольна робота \No 2, варіант  \No \textbf { 67 }\par}}
{
%\begin {large}

Тестова частина. Завдання 1-5: Що буде виведено за виконання (повідомлення інтерпретатора про помилку позначати error)?

Задачі (на звороті): написати функцію для розв'язання задачі.

\vspace{5mm} \begin{minipage}[t]{8.5cm}
               
{\begin{center}Завдання {1}\end{center}}
\vspace{-6mm}

\begin{verbatim}
a,b,c=8,1,2
def g(b):
    a=3
    b=4
    c=1
    return a+b+c

a,b,c=6,4,7
print(g(b),a,b,c)
\end{verbatim}


\vspace{5mm}
               
{\begin{center}Завдання {2}\end{center}}
\vspace{-6mm}

\begin{verbatim}
def f():
    try:
        res = int(2//0)
    except BaseException: return 7
    except TypeError: return 1
    return res

print(f())
\end{verbatim}


\vspace{5mm}
               
{\begin{center}Завдання {3}\end{center}}
\vspace{-6mm}

\begin{verbatim}
n=9
while n>=3:
    n+=-3
print(n)
\end{verbatim}


\end{minipage}
\vline \hspace{6mm}
\begin{minipage}[t]{10.7cm}
               
{\begin{center}Завдання {4}\end{center}}
\vspace{-6mm}

\begin{verbatim}
for d in range(-3,-3+6,-2):
    if d>=3:
        break
    print(d, end=' ')
    d=7
else:
    print(d, end=' ')
print(d, end=' ')
\end{verbatim}


\vspace{5mm}
               
{\begin{center}Завдання {5}\end{center}}
\vspace{-6mm}

\begin{verbatim}
def f(n):
    if n > 9:
        f(n // 10)
    else:
        print(n % 10, end=' ')

f(543216)
\end{verbatim}


\end{minipage}

\newpage

{{22.11.2024 Контрольна робота \No 2, варіант  \No \textbf { 68 }\par}}
{
%\begin {large}

Тестова частина. Завдання 1-5: Що буде виведено за виконання (повідомлення інтерпретатора про помилку позначати error)?

Задачі (на звороті): написати функцію для розв'язання задачі.

\vspace{5mm} \begin{minipage}[t]{8.5cm}
               
{\begin{center}Завдання {1}\end{center}}
\vspace{-6mm}

\begin{verbatim}
a,b,c=3,9,4
def g(a):
    global c
    a+=3
    b=1
    c=4
    return a+b+c

a,b,c=1,5,7
print(g(b),a,b,c)
\end{verbatim}


\vspace{5mm}
               
{\begin{center}Завдання {2}\end{center}}
\vspace{-6mm}

\begin{verbatim}
def f():
    try:
        res = int("b9")
        return 40
    except ValueError: return 8
    except ZeroDivisionError: return 4
    else: return 30
    finally: return 21
    return res

print(f())
\end{verbatim}


\vspace{5mm}
               
{\begin{center}Завдання {3}\end{center}}
\vspace{-6mm}

\begin{verbatim}
n=5
while n<=5:
    n+=1
print(n)
\end{verbatim}


\end{minipage}
\vline \hspace{6mm}
\begin{minipage}[t]{10.7cm}
               
{\begin{center}Завдання {4}\end{center}}
\vspace{-6mm}

\begin{verbatim}
for a in range(0,0+4,2):
    if a<8:
        pass
    print(a, end=' ')
    a=-5
else:
    print(a, end=' ')
print(a, end=' ')
\end{verbatim}


\vspace{5mm}
               
{\begin{center}Завдання {5}\end{center}}
\vspace{-6mm}

\begin{verbatim}
def f(n):
    print(n % 10, end=' ')
    if n > 9:
        f(n // 10)
    

f(543216)
\end{verbatim}


\end{minipage}

\newpage

{{22.11.2024 Контрольна робота \No 2, варіант  \No \textbf { 69 }\par}}
{
%\begin {large}

Тестова частина. Завдання 1-5: Що буде виведено за виконання (повідомлення інтерпретатора про помилку позначати error)?

Задачі (на звороті): написати функцію для розв'язання задачі.

\vspace{5mm} \begin{minipage}[t]{8.5cm}
               
{\begin{center}Завдання {1}\end{center}}
\vspace{-6mm}

\begin{verbatim}
a,b,c=8,6,0
def h(b):
    a=2
    b-=5
    c=5
    return a+b+c

a,b,c=2,4,8
print(h(b),a,b,c)
\end{verbatim}


\vspace{5mm}
               
{\begin{center}Завдання {2}\end{center}}
\vspace{-6mm}

\begin{verbatim}
def f():
    try:
        res = 5!=0
    except KeyboardInterrupt: return 6
    except Exception: return 7
    finally: return 24
    return res

print(f())
\end{verbatim}


\vspace{5mm}
               
{\begin{center}Завдання {3}\end{center}}
\vspace{-6mm}

\begin{verbatim}
t=4
while t<9:
    t+=2
print(t)
\end{verbatim}


\end{minipage}
\vline \hspace{6mm}
\begin{minipage}[t]{10.7cm}
               
{\begin{center}Завдання {4}\end{center}}
\vspace{-6mm}

\begin{verbatim}
for c in range(-7,-7+5,2):
    if c<=5:
        break
    print(c, end=' ')
    c=8
    if c>=6:
        break
else:
    print(c, end=' ')
print(c, end=' ')
\end{verbatim}


\vspace{5mm}
               
{\begin{center}Завдання {5}\end{center}}
\vspace{-6mm}

\begin{verbatim}
def f(n):
    print(n % 10, end=' ')
    if n > 9:
           f(n // 100)
    

f(123456)
\end{verbatim}


\end{minipage}

\newpage

{{22.11.2024 Контрольна робота \No 2, варіант  \No \textbf { 70 }\par}}
{
%\begin {large}

Тестова частина. Завдання 1-5: Що буде виведено за виконання (повідомлення інтерпретатора про помилку позначати error)?

Задачі (на звороті): написати функцію для розв'язання задачі.

\vspace{5mm} \begin{minipage}[t]{8.5cm}
               
{\begin{center}Завдання {1}\end{center}}
\vspace{-6mm}

\begin{verbatim}
a,b,c=2,6,5
def g(b):
    global c
    a=4
    b*=3
    c=1
    return a+b+c

a,b,c=7,9,3
print(g(a),a,b,c)
\end{verbatim}


\vspace{5mm}
               
{\begin{center}Завдання {2}\end{center}}
\vspace{-6mm}

\begin{verbatim}
def f():
    try:
        res = int("0")
        return 41
    except BaseException: return 2
    except Exception: return 3
    else: return 32
    return res

print(f())
\end{verbatim}


\vspace{5mm}
               
{\begin{center}Завдання {3}\end{center}}
\vspace{-6mm}

\begin{verbatim}
j=0
while j<=11:
    j+=1
print(j)
\end{verbatim}


\end{minipage}
\vline \hspace{6mm}
\begin{minipage}[t]{10.7cm}
               
{\begin{center}Завдання {4}\end{center}}
\vspace{-6mm}

\begin{verbatim}
for e in range(6,6+3,-3):
    if e<7:
        continue
    print(e, end=' ')
    e=-6
else:
    print(13, end=' ')
print(e, end=' ')
\end{verbatim}


\vspace{5mm}
               
{\begin{center}Завдання {5}\end{center}}
\vspace{-6mm}

\begin{verbatim}
def f(n):
    print(n % 10, end=' ')
    if n > 9:
        f(n // 100)
    

f(987654)
\end{verbatim}


\end{minipage}

\newpage

{{22.11.2024 Контрольна робота \No 2, варіант  \No \textbf { 71 }\par}}
{
%\begin {large}

Тестова частина. Завдання 1-5: Що буде виведено за виконання (повідомлення інтерпретатора про помилку позначати error)?

Задачі (на звороті): написати функцію для розв'язання задачі.

\vspace{5mm} \begin{minipage}[t]{8.5cm}
               
{\begin{center}Завдання {1}\end{center}}
\vspace{-6mm}

\begin{verbatim}
a,b,c=3,2,8
def f(a):
    global c
    a=2
    b=3
    c=5
    return a+b+c

a,b,c=9,6,1
print(f(b),a,b,c)
\end{verbatim}


\vspace{5mm}
               
{\begin{center}Завдання {2}\end{center}}
\vspace{-6mm}

\begin{verbatim}
def f():
    try:
        res = int(4/0.0)
        return 42
    except ZeroDivisionError: return 9
    except TypeError: return 7
    else: return 31
    return res

print(f())
\end{verbatim}


\vspace{5mm}
               
{\begin{center}Завдання {3}\end{center}}
\vspace{-6mm}

\begin{verbatim}
k=8
while k<9:
    k+=2
print(k)
\end{verbatim}


\end{minipage}
\vline \hspace{6mm}
\begin{minipage}[t]{10.7cm}
               
{\begin{center}Завдання {4}\end{center}}
\vspace{-6mm}

\begin{verbatim}
for b in range(-10,-10+6,-2):
    if b>=6:
        pass
    print(b, end=' ')
    b=-7
else:
    print(b, end=' ')
print(b, end=' ')
\end{verbatim}


\vspace{5mm}
               
{\begin{center}Завдання {5}\end{center}}
\vspace{-6mm}

\begin{verbatim}
def f(n):
    if n > 9:
        f(n // 100)
    print(n % 10, end=' ')
    

f(123456)
\end{verbatim}


\end{minipage}

\newpage

{{22.11.2024 Контрольна робота \No 2, варіант  \No \textbf { 72 }\par}}
{
%\begin {large}

Тестова частина. Завдання 1-5: Що буде виведено за виконання (повідомлення інтерпретатора про помилку позначати error)?

Задачі (на звороті): написати функцію для розв'язання задачі.

\vspace{5mm} \begin{minipage}[t]{8.5cm}
               
{\begin{center}Завдання {1}\end{center}}
\vspace{-6mm}

\begin{verbatim}
a,b,c=1,0,6
def f(a):
    a=2
    b+=1
    c=3
    return a+b+c

a,b,c=2,1,7
print(f(b),a,b,c)
\end{verbatim}


\vspace{5mm}
               
{\begin{center}Завдання {2}\end{center}}
\vspace{-6mm}

\begin{verbatim}
def f():
    try:
        res = int(5%2)
    except KeyboardInterrupt: return 6
    except ValueError: return 2
    finally: return 22
    return res

print(f())
\end{verbatim}


\vspace{5mm}
               
{\begin{center}Завдання {3}\end{center}}
\vspace{-6mm}

\begin{verbatim}
j=10
while j>8:
    j+=-2
print(j)
\end{verbatim}


\end{minipage}
\vline \hspace{6mm}
\begin{minipage}[t]{10.7cm}
               
{\begin{center}Завдання {4}\end{center}}
\vspace{-6mm}

\begin{verbatim}
for f in range(-1,-1+2,3):
    if f>4:
        continue
    print(f, end=' ')
    f=1
else:
    print(f, end=' ')
print(f, end=' ')
\end{verbatim}


\vspace{5mm}
               
{\begin{center}Завдання {5}\end{center}}
\vspace{-6mm}

\begin{verbatim}
def f(n):
    if n > 2:
        f(n // 2)
    else:
        print(n % 2, end=' ')

f(16)
\end{verbatim}


\end{minipage}

\newpage

{{22.11.2024 Контрольна робота \No 2, варіант  \No \textbf { 73 }\par}}
{
%\begin {large}

Тестова частина. Завдання 1-5: Що буде виведено за виконання (повідомлення інтерпретатора про помилку позначати error)?

Задачі (на звороті): написати функцію для розв'язання задачі.

\vspace{5mm} \begin{minipage}[t]{8.5cm}
               
{\begin{center}Завдання {1}\end{center}}
\vspace{-6mm}

\begin{verbatim}
a,b,c=9,3,8
def g(b):
    global c
    a=4
    b-=2
    c=5
    return a+b+c

a,b,c=5,4,0
print(g(a),a,b,c)
\end{verbatim}


\vspace{5mm}
               
{\begin{center}Завдання {2}\end{center}}
\vspace{-6mm}

\begin{verbatim}
def f():
    try:
        res = int("3")
    except ZeroDivisionError: return 1
    except Exception: return 8
    finally: return 21
    return res

print(f())
\end{verbatim}


\vspace{5mm}
               
{\begin{center}Завдання {3}\end{center}}
\vspace{-6mm}

\begin{verbatim}
m=11
while m>=0:
    m+=-3
print(m)
\end{verbatim}


\end{minipage}
\vline \hspace{6mm}
\begin{minipage}[t]{10.7cm}
               
{\begin{center}Завдання {4}\end{center}}
\vspace{-6mm}

\begin{verbatim}
for f in range(9,9+5,-2):
    if f>8:
        continue
    print(f, end=' ')
    f=-10
else:
    print(f, end=' ')
print(f, end=' ')
\end{verbatim}


\vspace{5mm}
               
{\begin{center}Завдання {5}\end{center}}
\vspace{-6mm}

\begin{verbatim}
def f(n):
    if n > 9:
        return f(n // 10) + 1
    else:
        return 1

print(f(123456))
\end{verbatim}


\end{minipage}

\newpage

{{22.11.2024 Контрольна робота \No 2, варіант  \No \textbf { 74 }\par}}
{
%\begin {large}

Тестова частина. Завдання 1-5: Що буде виведено за виконання (повідомлення інтерпретатора про помилку позначати error)?

Задачі (на звороті): написати функцію для розв'язання задачі.

\vspace{5mm} \begin{minipage}[t]{8.5cm}
               
{\begin{center}Завдання {1}\end{center}}
\vspace{-6mm}

\begin{verbatim}
a,b,c=0,3,7
def h(a):
    a=4
    b+=2
    c=5
    return a+b+c

a,b,c=4,8,2
print(h(b),a,b,c)
\end{verbatim}


\vspace{5mm}
               
{\begin{center}Завдання {2}\end{center}}
\vspace{-6mm}

\begin{verbatim}
def f():
    try:
        res = int("a0")
        return 42
    except ValueError: return 2
    except TypeError: return 8
    else: return 35
    return res

print(f())
\end{verbatim}


\vspace{5mm}
               
{\begin{center}Завдання {3}\end{center}}
\vspace{-6mm}

\begin{verbatim}
i=1
while i<7:
    i+=1
print(i)
\end{verbatim}


\end{minipage}
\vline \hspace{6mm}
\begin{minipage}[t]{10.7cm}
               
{\begin{center}Завдання {4}\end{center}}
\vspace{-6mm}

\begin{verbatim}
for b in range(7,7+4,2):
    if b<5:
        pass
    print(b, end=' ')
    b=-8
else:
    print(b, end=' ')
print(b, end=' ')
\end{verbatim}


\vspace{5mm}
               
{\begin{center}Завдання {5}\end{center}}
\vspace{-6mm}

\begin{verbatim}
def f(n):
    if n > 9:
        f(n // 100)
    else:
        print(n % 10, end=' ')

f(987654)
\end{verbatim}


\end{minipage}

\newpage

{{22.11.2024 Контрольна робота \No 2, варіант  \No \textbf { 75 }\par}}
{
%\begin {large}

Тестова частина. Завдання 1-5: Що буде виведено за виконання (повідомлення інтерпретатора про помилку позначати error)?

Задачі (на звороті): написати функцію для розв'язання задачі.

\vspace{5mm} \begin{minipage}[t]{8.5cm}
               
{\begin{center}Завдання {1}\end{center}}
\vspace{-6mm}

\begin{verbatim}
a,b,c=5,1,9
def f(a):
    global c
    a*=3
    b=1
    c=5
    return a+b+c

a,b,c=6,9,1
print(f(b),a,b,c)
\end{verbatim}


\vspace{5mm}
               
{\begin{center}Завдання {2}\end{center}}
\vspace{-6mm}

\begin{verbatim}
def f():
    try:
        res = 6>=1
        return 45
    except KeyboardInterrupt: return 3
    except BaseException: return 4
    return res

print(f())
\end{verbatim}


\vspace{5mm}
               
{\begin{center}Завдання {3}\end{center}}
\vspace{-6mm}

\begin{verbatim}
m=3
while m<=10:
    m+=2
print(m)
\end{verbatim}


\end{minipage}
\vline \hspace{6mm}
\begin{minipage}[t]{10.7cm}
               
{\begin{center}Завдання {4}\end{center}}
\vspace{-6mm}

\begin{verbatim}
for c in range(-9,-9+3,3):
    if c<=6:
        continue
    print(c, end=' ')
    c=0
else:
    print(c, end=' ')
print(c, end=' ')
\end{verbatim}


\vspace{5mm}
               
{\begin{center}Завдання {5}\end{center}}
\vspace{-6mm}

\begin{verbatim}
def f(n):
    print(n % 2, end=' ')
    if n > 2:
        f(n // 2)
    

f(16)
\end{verbatim}


\end{minipage}

\newpage

{{22.11.2024 Контрольна робота \No 2, варіант  \No \textbf { 76 }\par}}
{
%\begin {large}

Тестова частина. Завдання 1-5: Що буде виведено за виконання (повідомлення інтерпретатора про помилку позначати error)?

Задачі (на звороті): написати функцію для розв'язання задачі.

\vspace{5mm} \begin{minipage}[t]{8.5cm}
               
{\begin{center}Завдання {1}\end{center}}
\vspace{-6mm}

\begin{verbatim}
a,b,c=6,4,8
def f(b):
    a-=1
    b=2
    c=4
    return a+b+c

a,b,c=0,5,2
print(f(a),a,b,c)
\end{verbatim}


\vspace{5mm}
               
{\begin{center}Завдання {2}\end{center}}
\vspace{-6mm}

\begin{verbatim}
def f():
    try:
        res = 5<3
    except ZeroDivisionError: return 1
    except KeyboardInterrupt: return 0
    else: return 31
    finally: return 20
    return res

print(f())
\end{verbatim}


\vspace{5mm}
               
{\begin{center}Завдання {3}\end{center}}
\vspace{-6mm}

\begin{verbatim}
t=8
while t<8:
    t+=2
print(t)
\end{verbatim}


\end{minipage}
\vline \hspace{6mm}
\begin{minipage}[t]{10.7cm}
               
{\begin{center}Завдання {4}\end{center}}
\vspace{-6mm}

\begin{verbatim}
for a in range(-5,-5+6,-3):
    if a>=4:
        break
    print(a, end=' ')
    a=-4
else:
    print(a, end=' ')
print(a, end=' ')
\end{verbatim}


\vspace{5mm}
               
{\begin{center}Завдання {5}\end{center}}
\vspace{-6mm}

\begin{verbatim}
def f(n):
    if n <= 9:
        print(n % 10, end=' ')
    else:
        f(n // 10)

f(123456)
\end{verbatim}


\end{minipage}

\newpage

{{22.11.2024 Контрольна робота \No 2, варіант  \No \textbf { 77 }\par}}
{
%\begin {large}

Тестова частина. Завдання 1-5: Що буде виведено за виконання (повідомлення інтерпретатора про помилку позначати error)?

Задачі (на звороті): написати функцію для розв'язання задачі.

\vspace{5mm} \begin{minipage}[t]{8.5cm}
               
{\begin{center}Завдання {1}\end{center}}
\vspace{-6mm}

\begin{verbatim}
a,b,c=3,7,7
def h(a):
    global c
    a=3
    b+=5
    c=4
    return a+b+c

a,b,c=6,5,1
print(h(a),a,b,c)
\end{verbatim}


\vspace{5mm}
               
{\begin{center}Завдання {2}\end{center}}
\vspace{-6mm}

\begin{verbatim}
def f():
    try:
        res = int(7%1)
        return 44
    except Exception: return 9
    except ValueError: return 6
    else: return 30
    return res

print(f())
\end{verbatim}


\vspace{5mm}
               
{\begin{center}Завдання {3}\end{center}}
\vspace{-6mm}

\begin{verbatim}
k=5
while k>4:
    k+=-2
print(k)
\end{verbatim}


\end{minipage}
\vline \hspace{6mm}
\begin{minipage}[t]{10.7cm}
               
{\begin{center}Завдання {4}\end{center}}
\vspace{-6mm}

\begin{verbatim}
for d in range(2,2+2,3):
    if d>3:
        continue
    print(d, end=' ')
    d=3
else:
    print(13, end=' ')
print(d, end=' ')
\end{verbatim}


\vspace{5mm}
               
{\begin{center}Завдання {5}\end{center}}
\vspace{-6mm}

\begin{verbatim}
def f(n):
    print(n % 10, end=' ')
    if n > 9:
        f(n // 10)
    

f(123456)
\end{verbatim}


\end{minipage}

\newpage

{{22.11.2024 Контрольна робота \No 2, варіант  \No \textbf { 78 }\par}}
{
%\begin {large}

Тестова частина. Завдання 1-5: Що буде виведено за виконання (повідомлення інтерпретатора про помилку позначати error)?

Задачі (на звороті): написати функцію для розв'язання задачі.

\vspace{5mm} \begin{minipage}[t]{8.5cm}
               
{\begin{center}Завдання {1}\end{center}}
\vspace{-6mm}

\begin{verbatim}
a,b,c=0,5,8
def g(a):
    global c
    a=4
    b=3
    c=2
    return a+b+c

a,b,c=3,6,0
print(g(a),a,b,c)
\end{verbatim}


\vspace{5mm}
               
{\begin{center}Завдання {2}\end{center}}
\vspace{-6mm}

\begin{verbatim}
def f():
    try:
        res = int("0")
    except TypeError: return 2
    except BaseException: return 3
    finally: return 23
    return res

print(f())
\end{verbatim}


\vspace{5mm}
               
{\begin{center}Завдання {3}\end{center}}
\vspace{-6mm}

\begin{verbatim}
n=7
while n<=13:
    n+=1
print(n)
\end{verbatim}


\end{minipage}
\vline \hspace{6mm}
\begin{minipage}[t]{10.7cm}
               
{\begin{center}Завдання {4}\end{center}}
\vspace{-6mm}

\begin{verbatim}
for e in range(-6,-6+5,2):
    if e<7:
        break
    print(e, end=' ')
    e=5
else:
    print(e, end=' ')
print(e, end=' ')
\end{verbatim}


\vspace{5mm}
               
{\begin{center}Завдання {5}\end{center}}
\vspace{-6mm}

\begin{verbatim}
def f(n):
    if n <= 9:
        print(n % 10, end=' ')
    else:
        f(n // 10)

f(987651)
\end{verbatim}


\end{minipage}

\newpage

{{22.11.2024 Контрольна робота \No 2, варіант  \No \textbf { 79 }\par}}
{
%\begin {large}

Тестова частина. Завдання 1-5: Що буде виведено за виконання (повідомлення інтерпретатора про помилку позначати error)?

Задачі (на звороті): написати функцію для розв'язання задачі.

\vspace{5mm} \begin{minipage}[t]{8.5cm}
               
{\begin{center}Завдання {1}\end{center}}
\vspace{-6mm}

\begin{verbatim}
a,b,c=4,7,5
def f(b):
    a=4
    b=1
    c=5
    return a+b+c

a,b,c=9,2,1
print(f(a),a,b,c)
\end{verbatim}


\vspace{5mm}
               
{\begin{center}Завдання {2}\end{center}}
\vspace{-6mm}

\begin{verbatim}
def f():
    try:
        res = int(8//0)
    except ValueError: return 9
    except Exception: return 4
    else: return 33
    return res

print(f())
\end{verbatim}


\vspace{5mm}
               
{\begin{center}Завдання {3}\end{center}}
\vspace{-6mm}

\begin{verbatim}
i=6
while i>9:
    i+=-2
print(i)
\end{verbatim}


\end{minipage}
\vline \hspace{6mm}
\begin{minipage}[t]{10.7cm}
               
{\begin{center}Завдання {4}\end{center}}
\vspace{-6mm}

\begin{verbatim}
for f in range(-4,-4+3,-3):
    if f>=8:
        continue
    print(f, end=' ')
    f=10
    if f<=7:
        break
else:
    print(13, end=' ')
print(f, end=' ')
\end{verbatim}


\vspace{5mm}
               
{\begin{center}Завдання {5}\end{center}}
\vspace{-6mm}

\begin{verbatim}
def f(n):
    if n > 9:
        f(n // 10)
    else:
        print(n % 10, end=' ')

f(123456)
\end{verbatim}


\end{minipage}

\newpage

{{22.11.2024 Контрольна робота \No 2, варіант  \No \textbf { 80 }\par}}
{
%\begin {large}

Тестова частина. Завдання 1-5: Що буде виведено за виконання (повідомлення інтерпретатора про помилку позначати error)?

Задачі (на звороті): написати функцію для розв'язання задачі.

\vspace{5mm} \begin{minipage}[t]{8.5cm}
               
{\begin{center}Завдання {1}\end{center}}
\vspace{-6mm}

\begin{verbatim}
a,b,c=0,8,3
def g(b):
    a=3
    b*=1
    c=2
    return a+b+c

a,b,c=4,9,2
print(g(b),a,b,c)
\end{verbatim}


\vspace{5mm}
               
{\begin{center}Завдання {2}\end{center}}
\vspace{-6mm}

\begin{verbatim}
def f():
    try:
        res = int("d5")
        return 41
    except ZeroDivisionError: return 1
    except BaseException: return 7
    finally: return 22
    return res

print(f())
\end{verbatim}


\vspace{5mm}
               
{\begin{center}Завдання {3}\end{center}}
\vspace{-6mm}

\begin{verbatim}
t=5
while t<=12:
    t+=2
print(t)
\end{verbatim}


\end{minipage}
\vline \hspace{6mm}
\begin{minipage}[t]{10.7cm}
               
{\begin{center}Завдання {4}\end{center}}
\vspace{-6mm}

\begin{verbatim}
for e in range(8,8+4,-2):
    if e<=5:
        pass
    print(e, end=' ')
    e=4
else:
    print(e, end=' ')
print(e, end=' ')
\end{verbatim}


\vspace{5mm}
               
{\begin{center}Завдання {5}\end{center}}
\vspace{-6mm}

\begin{verbatim}
def f(n):
    if n > 9:
            f(n // 100)
    print(n % 10, end=' ')
    

f(987654)
\end{verbatim}


\end{minipage}

\newpage

{{22.11.2024 Контрольна робота \No 2, варіант  \No \textbf { 81 }\par}}
{
%\begin {large}

Тестова частина. Завдання 1-5: Що буде виведено за виконання (повідомлення інтерпретатора про помилку позначати error)?

Задачі (на звороті): написати функцію для розв'язання задачі.

\vspace{5mm} \begin{minipage}[t]{8.5cm}
               
{\begin{center}Завдання {1}\end{center}}
\vspace{-6mm}

\begin{verbatim}
a,b,c=7,1,5
def g(a):
    global c
    a=4
    b*=3
    c=5
    return a+b+c

a,b,c=0,9,8
print(g(a),a,b,c)
\end{verbatim}


\vspace{5mm}
               
{\begin{center}Завдання {2}\end{center}}
\vspace{-6mm}

\begin{verbatim}
def f():
    try:
        res = int("d9")
        return 40
    except KeyboardInterrupt: return 7
    except TypeError: return 2
    return res

print(f())
\end{verbatim}


\vspace{5mm}
               
{\begin{center}Завдання {3}\end{center}}
\vspace{-6mm}

\begin{verbatim}
t=12
while t>=5:
    t+=-3
print(t)
\end{verbatim}


\end{minipage}
\vline \hspace{6mm}
\begin{minipage}[t]{10.7cm}
               
{\begin{center}Завдання {4}\end{center}}
\vspace{-6mm}

\begin{verbatim}
for d in range(3,3+5,2):
    if d<4:
        break
    print(d, end=' ')
    d=6
else:
    print(d, end=' ')
print(d, end=' ')
\end{verbatim}


\vspace{5mm}
               
{\begin{center}Завдання {5}\end{center}}
\vspace{-6mm}

\begin{verbatim}
def f(n):
    if n > 9:
        f(n // 10)
    print(n % 10, end=' ')
    

f(123456)
\end{verbatim}


\end{minipage}

\newpage

{{22.11.2024 Контрольна робота \No 2, варіант  \No \textbf { 82 }\par}}
{
%\begin {large}

Тестова частина. Завдання 1-5: Що буде виведено за виконання (повідомлення інтерпретатора про помилку позначати error)?

Задачі (на звороті): написати функцію для розв'язання задачі.

\vspace{5mm} \begin{minipage}[t]{8.5cm}
               
{\begin{center}Завдання {1}\end{center}}
\vspace{-6mm}

\begin{verbatim}
a,b,c=4,6,2
def f(b):
    a=1
    b-=2
    c=3
    return a+b+c

a,b,c=3,5,7
print(f(b),a,b,c)
\end{verbatim}


\vspace{5mm}
               
{\begin{center}Завдання {2}\end{center}}
\vspace{-6mm}

\begin{verbatim}
def f():
    try:
        res = 0==4
    except ValueError: return 8
    except TypeError: return 3
    else: return 34
    finally: return 25
    return res

print(f())
\end{verbatim}


\vspace{5mm}
               
{\begin{center}Завдання {3}\end{center}}
\vspace{-6mm}

\begin{verbatim}
j=15
while j>=1:
    j+=-3
print(j)
\end{verbatim}


\end{minipage}
\vline \hspace{6mm}
\begin{minipage}[t]{10.7cm}
               
{\begin{center}Завдання {4}\end{center}}
\vspace{-6mm}

\begin{verbatim}
for a in range(4,4+3,-3):
    if a>=7:
        continue
    print(a, end=' ')
    a=9
else:
    print(a, end=' ')
print(a, end=' ')
\end{verbatim}


\vspace{5mm}
               
{\begin{center}Завдання {5}\end{center}}
\vspace{-6mm}

\begin{verbatim}
def f(n):
    if n > 9:
            f(n // 100)
    print(n % 10, end=' ')
    

f(987654)
\end{verbatim}


\end{minipage}

\newpage

{{22.11.2024 Контрольна робота \No 2, варіант  \No \textbf { 83 }\par}}
{
%\begin {large}

Тестова частина. Завдання 1-5: Що буде виведено за виконання (повідомлення інтерпретатора про помилку позначати error)?

Задачі (на звороті): написати функцію для розв'язання задачі.

\vspace{5mm} \begin{minipage}[t]{8.5cm}
               
{\begin{center}Завдання {1}\end{center}}
\vspace{-6mm}

\begin{verbatim}
a,b,c=7,5,3
def h(a):
    global c
    a=5
    b=4
    c=2
    return a+b+c

a,b,c=9,6,1
print(h(b),a,b,c)
\end{verbatim}


\vspace{5mm}
               
{\begin{center}Завдання {2}\end{center}}
\vspace{-6mm}

\begin{verbatim}
def f():
    try:
        res = int(5/0.0)
    except BaseException: return 4
    except Exception: return 1
    else: return 32
    finally: return 24
    return res

print(f())
\end{verbatim}


\vspace{5mm}
               
{\begin{center}Завдання {3}\end{center}}
\vspace{-6mm}

\begin{verbatim}
k=13
while k>6:
    k+=-2
print(k)
\end{verbatim}


\end{minipage}
\vline \hspace{6mm}
\begin{minipage}[t]{10.7cm}
               
{\begin{center}Завдання {4}\end{center}}
\vspace{-6mm}

\begin{verbatim}
for c in range(10,10+6,3):
    if c<=6:
        pass
    print(c, end=' ')
    c=-2
else:
    print(c, end=' ')
print(c, end=' ')
\end{verbatim}


\vspace{5mm}
               
{\begin{center}Завдання {5}\end{center}}
\vspace{-6mm}

\begin{verbatim}
def f(n):
    if n > 9:
        f(n // 10)
    else:
        print(n % 10, end=' ')

f(543216)
\end{verbatim}


\end{minipage}

\newpage

{{22.11.2024 Контрольна робота \No 2, варіант  \No \textbf { 84 }\par}}
{
%\begin {large}

Тестова частина. Завдання 1-5: Що буде виведено за виконання (повідомлення інтерпретатора про помилку позначати error)?

Задачі (на звороті): написати функцію для розв'язання задачі.

\vspace{5mm} \begin{minipage}[t]{8.5cm}
               
{\begin{center}Завдання {1}\end{center}}
\vspace{-6mm}

\begin{verbatim}
a,b,c=9,3,8
def h(a):
    global c
    a+=4
    b=1
    c=2
    return a+b+c

a,b,c=1,0,6
print(h(b),a,b,c)
\end{verbatim}


\vspace{5mm}
               
{\begin{center}Завдання {2}\end{center}}
\vspace{-6mm}

\begin{verbatim}
def f():
    try:
        res = int(6/1)
        return 43
    except ZeroDivisionError: return 8
    except KeyboardInterrupt: return 2
    return res

print(f())
\end{verbatim}


\vspace{5mm}
               
{\begin{center}Завдання {3}\end{center}}
\vspace{-6mm}

\begin{verbatim}
i=2
while i<6:
    i+=1
print(i)
\end{verbatim}


\end{minipage}
\vline \hspace{6mm}
\begin{minipage}[t]{10.7cm}
               
{\begin{center}Завдання {4}\end{center}}
\vspace{-6mm}

\begin{verbatim}
for b in range(-8,-8+4,-2):
    if b>3:
        continue
    print(b, end=' ')
    b=2
else:
    print(b, end=' ')
print(b, end=' ')
\end{verbatim}


\vspace{5mm}
               
{\begin{center}Завдання {5}\end{center}}
\vspace{-6mm}

\begin{verbatim}
def f(n):
    if n > 9:
        f(n // 10)
    print(n % 10, end=' ')
    

f(543216)
\end{verbatim}


\end{minipage}

\newpage

{{22.11.2024 Контрольна робота \No 2, варіант  \No \textbf { 85 }\par}}
{
%\begin {large}

Тестова частина. Завдання 1-5: Що буде виведено за виконання (повідомлення інтерпретатора про помилку позначати error)?

Задачі (на звороті): написати функцію для розв'язання задачі.

\vspace{5mm} \begin{minipage}[t]{8.5cm}
               
{\begin{center}Завдання {1}\end{center}}
\vspace{-6mm}

\begin{verbatim}
a,b,c=8,4,0
def f(b):
    a=1
    b=3
    c=1
    return a+b+c

a,b,c=2,7,2
print(f(b),a,b,c)
\end{verbatim}


\vspace{5mm}
               
{\begin{center}Завдання {2}\end{center}}
\vspace{-6mm}

\begin{verbatim}
def f():
    try:
        res = int("3")
        return 43
    except BaseException: return 4
    except TypeError: return 9
    finally: return 25
    return res

print(f())
\end{verbatim}


\vspace{5mm}
               
{\begin{center}Завдання {3}\end{center}}
\vspace{-6mm}

\begin{verbatim}
n=0
while n<=5:
    n+=2
print(n)
\end{verbatim}


\end{minipage}
\vline \hspace{6mm}
\begin{minipage}[t]{10.7cm}
               
{\begin{center}Завдання {4}\end{center}}
\vspace{-6mm}

\begin{verbatim}
for d in range(-5,-5+2,3):
    if d>5:
        break
    print(d, end=' ')
    d=4
    if d>5:
        break
else:
    print(d, end=' ')
print(d, end=' ')
\end{verbatim}


\vspace{5mm}
               
{\begin{center}Завдання {5}\end{center}}
\vspace{-6mm}

\begin{verbatim}
def f(n):
    print(n % 10, end=' ')
    if n > 9:
           f(n // 100)
    

f(123456)
\end{verbatim}


\end{minipage}

\newpage

{{22.11.2024 Контрольна робота \No 2, варіант  \No \textbf { 86 }\par}}
{
%\begin {large}

Тестова частина. Завдання 1-5: Що буде виведено за виконання (повідомлення інтерпретатора про помилку позначати error)?

Задачі (на звороті): написати функцію для розв'язання задачі.

\vspace{5mm} \begin{minipage}[t]{8.5cm}
               
{\begin{center}Завдання {1}\end{center}}
\vspace{-6mm}

\begin{verbatim}
a,b,c=4,2,0
def h(b):
    a+=5
    b=5
    c=4
    return a+b+c

a,b,c=2,4,6
print(h(b),a,b,c)
\end{verbatim}


\vspace{5mm}
               
{\begin{center}Завдання {2}\end{center}}
\vspace{-6mm}

\begin{verbatim}
def f():
    try:
        res = 0<=7
    except ValueError: return 1
    except ZeroDivisionError: return 6
    else: return 34
    return res

print(f())
\end{verbatim}


\vspace{5mm}
               
{\begin{center}Завдання {3}\end{center}}
\vspace{-6mm}

\begin{verbatim}
k=4
while k<10:
    k+=1
print(k)
\end{verbatim}


\end{minipage}
\vline \hspace{6mm}
\begin{minipage}[t]{10.7cm}
               
{\begin{center}Завдання {4}\end{center}}
\vspace{-6mm}

\begin{verbatim}
for b in range(7,7+2,-3):
    if b>=7:
        pass
    print(b, end=' ')
    b=-8
else:
    print(13, end=' ')
print(b, end=' ')
\end{verbatim}


\vspace{5mm}
               
{\begin{center}Завдання {5}\end{center}}
\vspace{-6mm}

\begin{verbatim}
def f(n):
    if n > 9:
        return f(n // 100) + 1
    else:
        return 1

print(f(123456))
\end{verbatim}


\end{minipage}

\newpage

{{22.11.2024 Контрольна робота \No 2, варіант  \No \textbf { 87 }\par}}
{
%\begin {large}

Тестова частина. Завдання 1-5: Що буде виведено за виконання (повідомлення інтерпретатора про помилку позначати error)?

Задачі (на звороті): написати функцію для розв'язання задачі.

\vspace{5mm} \begin{minipage}[t]{8.5cm}
               
{\begin{center}Завдання {1}\end{center}}
\vspace{-6mm}

\begin{verbatim}
a,b,c=5,9,7
def f(a):
    global c
    a-=2
    b=1
    c=3
    return a+b+c

a,b,c=1,3,8
print(f(a),a,b,c)
\end{verbatim}


\vspace{5mm}
               
{\begin{center}Завдання {2}\end{center}}
\vspace{-6mm}

\begin{verbatim}
def f():
    try:
        res = int(5//2)
        return 40
    except KeyboardInterrupt: return 7
    except Exception: return 7
    else: return 35
    return res

print(f())
\end{verbatim}


\vspace{5mm}
               
{\begin{center}Завдання {3}\end{center}}
\vspace{-6mm}

\begin{verbatim}
m=6
while m<9:
    m+=2
print(m)
\end{verbatim}


\end{minipage}
\vline \hspace{6mm}
\begin{minipage}[t]{10.7cm}
               
{\begin{center}Завдання {4}\end{center}}
\vspace{-6mm}

\begin{verbatim}
for a in range(4,4+5,-2):
    if a<6:
        continue
    print(a, end=' ')
    a=-10
else:
    print(a, end=' ')
print(a, end=' ')
\end{verbatim}


\vspace{5mm}
               
{\begin{center}Завдання {5}\end{center}}
\vspace{-6mm}

\begin{verbatim}
def f(n):
    print(n % 10, end=' ')
    if n > 9:
        f(n // 10)
    

f(123456)
\end{verbatim}


\end{minipage}

\newpage

{{22.11.2024 Контрольна робота \No 2, варіант  \No \textbf { 88 }\par}}
{
%\begin {large}

Тестова частина. Завдання 1-5: Що буде виведено за виконання (повідомлення інтерпретатора про помилку позначати error)?

Задачі (на звороті): написати функцію для розв'язання задачі.

\vspace{5mm} \begin{minipage}[t]{8.5cm}
               
{\begin{center}Завдання {1}\end{center}}
\vspace{-6mm}

\begin{verbatim}
a,b,c=5,8,0
def g(b):
    a=5
    b*=2
    c=3
    return a+b+c

a,b,c=3,6,4
print(g(a),a,b,c)
\end{verbatim}


\vspace{5mm}
               
{\begin{center}Завдання {2}\end{center}}
\vspace{-6mm}

\begin{verbatim}
def f():
    try:
        res = int("0")
    except TypeError: return 3
    except ZeroDivisionError: return 2
    finally: return 21
    return res

print(f())
\end{verbatim}


\vspace{5mm}
               
{\begin{center}Завдання {3}\end{center}}
\vspace{-6mm}

\begin{verbatim}
m=9
while m>=7:
    m+=-2
print(m)
\end{verbatim}


\end{minipage}
\vline \hspace{6mm}
\begin{minipage}[t]{10.7cm}
               
{\begin{center}Завдання {4}\end{center}}
\vspace{-6mm}

\begin{verbatim}
for c in range(10,10+6,2):
    if c<=4:
        continue
    print(c, end=' ')
    c=9
else:
    print(c, end=' ')
print(c, end=' ')
\end{verbatim}


\vspace{5mm}
               
{\begin{center}Завдання {5}\end{center}}
\vspace{-6mm}

\begin{verbatim}
def f(n):
    if n <= 9:
        print(n % 10, end=' ')
    else:
        f(n // 10)

f(987651)
\end{verbatim}


\end{minipage}

\newpage

{{22.11.2024 Контрольна робота \No 2, варіант  \No \textbf { 89 }\par}}
{
%\begin {large}

Тестова частина. Завдання 1-5: Що буде виведено за виконання (повідомлення інтерпретатора про помилку позначати error)?

Задачі (на звороті): написати функцію для розв'язання задачі.

\vspace{5mm} \begin{minipage}[t]{8.5cm}
               
{\begin{center}Завдання {1}\end{center}}
\vspace{-6mm}

\begin{verbatim}
a,b,c=7,1,2
def g(a):
    global c
    a-=4
    b=1
    c=2
    return a+b+c

a,b,c=9,7,2
print(g(b),a,b,c)
\end{verbatim}


\vspace{5mm}
               
{\begin{center}Завдання {2}\end{center}}
\vspace{-6mm}

\begin{verbatim}
def f():
    try:
        res = int(6%0)
    except Exception: return 8
    except KeyboardInterrupt: return 5
    finally: return 22
    return res

print(f())
\end{verbatim}


\vspace{5mm}
               
{\begin{center}Завдання {3}\end{center}}
\vspace{-6mm}

\begin{verbatim}
n=8
while n>9:
    n+=-3
print(n)
\end{verbatim}


\end{minipage}
\vline \hspace{6mm}
\begin{minipage}[t]{10.7cm}
               
{\begin{center}Завдання {4}\end{center}}
\vspace{-6mm}

\begin{verbatim}
for e in range(-9,-9+4,-3):
    if e>8:
        break
    print(e, end=' ')
    e=-3
else:
    print(e, end=' ')
print(e, end=' ')
\end{verbatim}


\vspace{5mm}
               
{\begin{center}Завдання {5}\end{center}}
\vspace{-6mm}

\begin{verbatim}
def f(n):
    if n > 9:
        f(n // 10)
    else:
        print(n % 10, end=' ')

f(123456)
\end{verbatim}


\end{minipage}

\newpage

{{22.11.2024 Контрольна робота \No 2, варіант  \No \textbf { 90 }\par}}
{
%\begin {large}

Тестова частина. Завдання 1-5: Що буде виведено за виконання (повідомлення інтерпретатора про помилку позначати error)?

Задачі (на звороті): написати функцію для розв'язання задачі.

\vspace{5mm} \begin{minipage}[t]{8.5cm}
               
{\begin{center}Завдання {1}\end{center}}
\vspace{-6mm}

\begin{verbatim}
a,b,c=3,5,9
def g(a):
    global c
    a=2
    b=3
    c=4
    return a+b+c

a,b,c=4,1,8
print(g(a),a,b,c)
\end{verbatim}


\vspace{5mm}
               
{\begin{center}Завдання {2}\end{center}}
\vspace{-6mm}

\begin{verbatim}
def f():
    try:
        res = int("b4")
        return 45
    except BaseException: return 1
    except ValueError: return 9
    else: return 33
    return res

print(f())
\end{verbatim}


\vspace{5mm}
               
{\begin{center}Завдання {3}\end{center}}
\vspace{-6mm}

\begin{verbatim}
j=3
while j<=6:
    j+=1
print(j)
\end{verbatim}


\end{minipage}
\vline \hspace{6mm}
\begin{minipage}[t]{10.7cm}
               
{\begin{center}Завдання {4}\end{center}}
\vspace{-6mm}

\begin{verbatim}
for f in range(5,5+3,2):
    if f>=3:
        continue
    print(f, end=' ')
    f=3
else:
    print(f, end=' ')
print(f, end=' ')
\end{verbatim}


\vspace{5mm}
               
{\begin{center}Завдання {5}\end{center}}
\vspace{-6mm}

\begin{verbatim}
def f(n):
    if n > 9:
        f(n // 100)
    print(n % 10, end=' ')
    

f(123456)
\end{verbatim}


\end{minipage}

\newpage

{{22.11.2024 Контрольна робота \No 2, варіант  \No \textbf { 91 }\par}}
{
%\begin {large}

Тестова частина. Завдання 1-5: Що буде виведено за виконання (повідомлення інтерпретатора про помилку позначати error)?

Задачі (на звороті): написати функцію для розв'язання задачі.

\vspace{5mm} \begin{minipage}[t]{8.5cm}
               
{\begin{center}Завдання {1}\end{center}}
\vspace{-6mm}

\begin{verbatim}
a,b,c=1,3,6
def f(b):
    a+=4
    b=3
    c=1
    return a+b+c

a,b,c=8,9,4
print(f(a),a,b,c)
\end{verbatim}


\vspace{5mm}
               
{\begin{center}Завдання {2}\end{center}}
\vspace{-6mm}

\begin{verbatim}
def f():
    try:
        res = int("c3")
        return 41
    except KeyboardInterrupt: return 4
    except ValueError: return 2
    else: return 32
    finally: return 24
    return res

print(f())
\end{verbatim}


\vspace{5mm}
               
{\begin{center}Завдання {3}\end{center}}
\vspace{-6mm}

\begin{verbatim}
t=1
while t<7:
    t+=1
print(t)
\end{verbatim}


\end{minipage}
\vline \hspace{6mm}
\begin{minipage}[t]{10.7cm}
               
{\begin{center}Завдання {4}\end{center}}
\vspace{-6mm}

\begin{verbatim}
for c in range(-7,-7+2,3):
    if c<=8:
        continue
    print(c, end=' ')
    c=2
else:
    print(c, end=' ')
print(c, end=' ')
\end{verbatim}


\vspace{5mm}
               
{\begin{center}Завдання {5}\end{center}}
\vspace{-6mm}

\begin{verbatim}
def f(n):
    if n <= 9:
        print(n % 10, end=' ')
    else:
        f(n // 10)

f(123456)
\end{verbatim}


\end{minipage}

\newpage

{{22.11.2024 Контрольна робота \No 2, варіант  \No \textbf { 92 }\par}}
{
%\begin {large}

Тестова частина. Завдання 1-5: Що буде виведено за виконання (повідомлення інтерпретатора про помилку позначати error)?

Задачі (на звороті): написати функцію для розв'язання задачі.

\vspace{5mm} \begin{minipage}[t]{8.5cm}
               
{\begin{center}Завдання {1}\end{center}}
\vspace{-6mm}

\begin{verbatim}
a,b,c=0,5,0
def h(a):
    a*=5
    b=3
    c=2
    return a+b+c

a,b,c=1,4,3
print(h(b),a,b,c)
\end{verbatim}


\vspace{5mm}
               
{\begin{center}Завдання {2}\end{center}}
\vspace{-6mm}

\begin{verbatim}
def f():
    try:
        res = int(9//0.0)
    except BaseException: return 5
    except Exception: return 1
    return res

print(f())
\end{verbatim}


\vspace{5mm}
               
{\begin{center}Завдання {3}\end{center}}
\vspace{-6mm}

\begin{verbatim}
i=7
while i>0:
    i+=-3
print(i)
\end{verbatim}


\end{minipage}
\vline \hspace{6mm}
\begin{minipage}[t]{10.7cm}
               
{\begin{center}Завдання {4}\end{center}}
\vspace{-6mm}

\begin{verbatim}
for e in range(6,6+5,-2):
    if e<5:
        pass
    print(e, end=' ')
    e=-9
    if e<3:
        break
else:
    print(e, end=' ')
print(e, end=' ')
\end{verbatim}


\vspace{5mm}
               
{\begin{center}Завдання {5}\end{center}}
\vspace{-6mm}

\begin{verbatim}
def f(n):
    if n > 2:
        f(n // 2)
    print(n % 2, end=' ')
    

f(16)
\end{verbatim}


\end{minipage}

\newpage

{{22.11.2024 Контрольна робота \No 2, варіант  \No \textbf { 93 }\par}}
{
%\begin {large}

Тестова частина. Завдання 1-5: Що буде виведено за виконання (повідомлення інтерпретатора про помилку позначати error)?

Задачі (на звороті): написати функцію для розв'язання задачі.

\vspace{5mm} \begin{minipage}[t]{8.5cm}
               
{\begin{center}Завдання {1}\end{center}}
\vspace{-6mm}

\begin{verbatim}
a,b,c=7,8,9
def f(b):
    global c
    a=5
    b-=4
    c=1
    return a+b+c

a,b,c=5,6,2
print(f(a),a,b,c)
\end{verbatim}


\vspace{5mm}
               
{\begin{center}Завдання {2}\end{center}}
\vspace{-6mm}

\begin{verbatim}
def f():
    try:
        res = int(0/2)
        return 44
    except ZeroDivisionError: return 8
    except TypeError: return 7
    else: return 31
    finally: return 23
    return res

print(f())
\end{verbatim}


\vspace{5mm}
               
{\begin{center}Завдання {3}\end{center}}
\vspace{-6mm}

\begin{verbatim}
i=9
while i<=12:
    i+=2
print(i)
\end{verbatim}


\end{minipage}
\vline \hspace{6mm}
\begin{minipage}[t]{10.7cm}
               
{\begin{center}Завдання {4}\end{center}}
\vspace{-6mm}

\begin{verbatim}
for d in range(3,3+3,-3):
    if d>=3:
        break
    print(d, end=' ')
    d=-6
else:
    print(d, end=' ')
print(d, end=' ')
\end{verbatim}


\vspace{5mm}
               
{\begin{center}Завдання {5}\end{center}}
\vspace{-6mm}

\begin{verbatim}
def f(n):
    if n > 9:
        f(n // 100)
    else:
        print(n % 10, end=' ')

f(123456)
\end{verbatim}


\end{minipage}

\newpage

{{22.11.2024 Контрольна робота \No 2, варіант  \No \textbf { 94 }\par}}
{
%\begin {large}

Тестова частина. Завдання 1-5: Що буде виведено за виконання (повідомлення інтерпретатора про помилку позначати error)?

Задачі (на звороті): написати функцію для розв'язання задачі.

\vspace{5mm} \begin{minipage}[t]{8.5cm}
               
{\begin{center}Завдання {1}\end{center}}
\vspace{-6mm}

\begin{verbatim}
a,b,c=0,6,6
def h(b):
    a=5
    b=1
    c=2
    return a+b+c

a,b,c=7,8,3
print(h(b),a,b,c)
\end{verbatim}


\vspace{5mm}
               
{\begin{center}Завдання {2}\end{center}}
\vspace{-6mm}

\begin{verbatim}
def f():
    try:
        res = 6<=5
    except KeyboardInterrupt: return 9
    except Exception: return 4
    return res

print(f())
\end{verbatim}


\vspace{5mm}
               
{\begin{center}Завдання {3}\end{center}}
\vspace{-6mm}

\begin{verbatim}
j=14
while j>=2:
    j+=-2
print(j)
\end{verbatim}


\end{minipage}
\vline \hspace{6mm}
\begin{minipage}[t]{10.7cm}
               
{\begin{center}Завдання {4}\end{center}}
\vspace{-6mm}

\begin{verbatim}
for a in range(-1,-1+6,2):
    if a<4:
        continue
    print(a, end=' ')
    a=-5
else:
    print(a, end=' ')
print(a, end=' ')
\end{verbatim}


\vspace{5mm}
               
{\begin{center}Завдання {5}\end{center}}
\vspace{-6mm}

\begin{verbatim}
def f(n):
    print(n % 10, end=' ')
    if n > 9:
        f(n // 100)
    

f(987654)
\end{verbatim}


\end{minipage}

\newpage

{{22.11.2024 Контрольна робота \No 2, варіант  \No \textbf { 95 }\par}}
{
%\begin {large}

Тестова частина. Завдання 1-5: Що буде виведено за виконання (повідомлення інтерпретатора про помилку позначати error)?

Задачі (на звороті): написати функцію для розв'язання задачі.

\vspace{5mm} \begin{minipage}[t]{8.5cm}
               
{\begin{center}Завдання {1}\end{center}}
\vspace{-6mm}

\begin{verbatim}
a,b,c=4,9,0
def h(b):
    global c
    a=5
    b=4
    c=3
    return a+b+c

a,b,c=2,5,1
print(h(b),a,b,c)
\end{verbatim}


\vspace{5mm}
               
{\begin{center}Завдання {2}\end{center}}
\vspace{-6mm}

\begin{verbatim}
def f():
    try:
        res = int("6")
        return 42
    except ZeroDivisionError: return 5
    except TypeError: return 8
    else: return 30
    return res

print(f())
\end{verbatim}


\vspace{5mm}
               
{\begin{center}Завдання {3}\end{center}}
\vspace{-6mm}

\begin{verbatim}
t=10
while t>1:
    t+=-3
print(t)
\end{verbatim}


\end{minipage}
\vline \hspace{6mm}
\begin{minipage}[t]{10.7cm}
               
{\begin{center}Завдання {4}\end{center}}
\vspace{-6mm}

\begin{verbatim}
for b in range(-2,-2+4,-2):
    if b>7:
        pass
    print(b, end=' ')
    b=-2
else:
    print(b, end=' ')
print(b, end=' ')
\end{verbatim}


\vspace{5mm}
               
{\begin{center}Завдання {5}\end{center}}
\vspace{-6mm}

\begin{verbatim}
def f(n):
    if n > 2:
        f(n // 2)
    else:
        print(n % 2, end=' ')

f(16)
\end{verbatim}


\end{minipage}

\newpage

{{22.11.2024 Контрольна робота \No 2, варіант  \No \textbf { 96 }\par}}
{
%\begin {large}

Тестова частина. Завдання 1-5: Що буде виведено за виконання (повідомлення інтерпретатора про помилку позначати error)?

Задачі (на звороті): написати функцію для розв'язання задачі.

\vspace{5mm} \begin{minipage}[t]{8.5cm}
               
{\begin{center}Завдання {1}\end{center}}
\vspace{-6mm}

\begin{verbatim}
a,b,c=1,8,6
def g(b):
    global c
    a+=3
    b=5
    c=2
    return a+b+c

a,b,c=2,9,4
print(g(b),a,b,c)
\end{verbatim}


\vspace{5mm}
               
{\begin{center}Завдання {2}\end{center}}
\vspace{-6mm}

\begin{verbatim}
def f():
    try:
        res = int(3%0)
    except ValueError: return 0
    except BaseException: return 7
    finally: return 20
    return res

print(f())
\end{verbatim}


\vspace{5mm}
               
{\begin{center}Завдання {3}\end{center}}
\vspace{-6mm}

\begin{verbatim}
n=6
while n<=11:
    n+=1
print(n)
\end{verbatim}


\end{minipage}
\vline \hspace{6mm}
\begin{minipage}[t]{10.7cm}
               
{\begin{center}Завдання {4}\end{center}}
\vspace{-6mm}

\begin{verbatim}
for f in range(-3,-3+4,3):
    if f<=6:
        continue
    print(f, end=' ')
    f=6
else:
    print(13, end=' ')
print(f, end=' ')
\end{verbatim}


\vspace{5mm}
               
{\begin{center}Завдання {5}\end{center}}
\vspace{-6mm}

\begin{verbatim}
def f(n):
    if n <= 9:
        print(n % 10, end=' ')
    else:
        f(n // 10)

f(543216)
\end{verbatim}


\end{minipage}

\newpage

{{22.11.2024 Контрольна робота \No 2, варіант  \No \textbf { 97 }\par}}
{
%\begin {large}

Тестова частина. Завдання 1-5: Що буде виведено за виконання (повідомлення інтерпретатора про помилку позначати error)?

Задачі (на звороті): написати функцію для розв'язання задачі.

\vspace{5mm} \begin{minipage}[t]{8.5cm}
               
{\begin{center}Завдання {1}\end{center}}
\vspace{-6mm}

\begin{verbatim}
a,b,c=0,3,7
def h(a):
    a*=4
    b=1
    c=4
    return a+b+c

a,b,c=5,6,1
print(h(a),a,b,c)
\end{verbatim}


\vspace{5mm}
               
{\begin{center}Завдання {2}\end{center}}
\vspace{-6mm}

\begin{verbatim}
def f():
    try:
        res = int("a2")
        return 44
    except TypeError: return 1
    except KeyboardInterrupt: return 7
    return res

print(f())
\end{verbatim}


\vspace{5mm}
               
{\begin{center}Завдання {3}\end{center}}
\vspace{-6mm}

\begin{verbatim}
k=7
while k>=6:
    k+=-2
print(k)
\end{verbatim}


\end{minipage}
\vline \hspace{6mm}
\begin{minipage}[t]{10.7cm}
               
{\begin{center}Завдання {4}\end{center}}
\vspace{-6mm}

\begin{verbatim}
for d in range(9,9+6,2):
    if d<=3:
        continue
    print(d, end=' ')
    d=8
else:
    print(d, end=' ')
print(d, end=' ')
\end{verbatim}


\vspace{5mm}
               
{\begin{center}Завдання {5}\end{center}}
\vspace{-6mm}

\begin{verbatim}
def f(n):
    if n > 9:
        return f(n // 10) + 1
    else:
        return 1

print(f(123456))
\end{verbatim}


\end{minipage}

\newpage

{{22.11.2024 Контрольна робота \No 2, варіант  \No \textbf { 98 }\par}}
{
%\begin {large}

Тестова частина. Завдання 1-5: Що буде виведено за виконання (повідомлення інтерпретатора про помилку позначати error)?

Задачі (на звороті): написати функцію для розв'язання задачі.

\vspace{5mm} \begin{minipage}[t]{8.5cm}
               
{\begin{center}Завдання {1}\end{center}}
\vspace{-6mm}

\begin{verbatim}
a,b,c=4,0,3
def f(a):
    a+=2
    b=3
    c=1
    return a+b+c

a,b,c=8,5,7
print(f(a),a,b,c)
\end{verbatim}


\vspace{5mm}
               
{\begin{center}Завдання {2}\end{center}}
\vspace{-6mm}

\begin{verbatim}
def f():
    try:
        res = int(0%1)
    except BaseException: return 8
    except ZeroDivisionError: return 6
    else: return 32
    finally: return 21
    return res

print(f())
\end{verbatim}


\vspace{5mm}
               
{\begin{center}Завдання {3}\end{center}}
\vspace{-6mm}

\begin{verbatim}
n=12
while n>8:
    n+=-3
print(n)
\end{verbatim}


\end{minipage}
\vline \hspace{6mm}
\begin{minipage}[t]{10.7cm}
               
{\begin{center}Завдання {4}\end{center}}
\vspace{-6mm}

\begin{verbatim}
for e in range(0,0+3,3):
    if e<7:
        pass
    print(e, end=' ')
    e=1
else:
    print(e, end=' ')
print(e, end=' ')
\end{verbatim}


\vspace{5mm}
               
{\begin{center}Завдання {5}\end{center}}
\vspace{-6mm}

\begin{verbatim}
def f(n):
    if n > 9:
        f(n // 100)
    else:
        print(n % 10, end=' ')

f(987654)
\end{verbatim}


\end{minipage}

\newpage

{{22.11.2024 Контрольна робота \No 2, варіант  \No \textbf { 99 }\par}}
{
%\begin {large}

Тестова частина. Завдання 1-5: Що буде виведено за виконання (повідомлення інтерпретатора про помилку позначати error)?

Задачі (на звороті): написати функцію для розв'язання задачі.

\vspace{5mm} \begin{minipage}[t]{8.5cm}
               
{\begin{center}Завдання {1}\end{center}}
\vspace{-6mm}

\begin{verbatim}
a,b,c=9,2,2
def g(b):
    global c
    a*=5
    b=5
    c=4
    return a+b+c

a,b,c=7,1,6
print(g(a),a,b,c)
\end{verbatim}


\vspace{5mm}
               
{\begin{center}Завдання {2}\end{center}}
\vspace{-6mm}

\begin{verbatim}
def f():
    try:
        res = 5<6
    except Exception: return 3
    except ValueError: return 9
    return res

print(f())
\end{verbatim}


\vspace{5mm}
               
{\begin{center}Завдання {3}\end{center}}
\vspace{-6mm}

\begin{verbatim}
m=14
while m>=3:
    m+=-2
print(m)
\end{verbatim}


\end{minipage}
\vline \hspace{6mm}
\begin{minipage}[t]{10.7cm}
               
{\begin{center}Завдання {4}\end{center}}
\vspace{-6mm}

\begin{verbatim}
for a in range(1,1+5,-3):
    if a>8:
        break
    print(a, end=' ')
    a=-1
else:
    print(13, end=' ')
print(a, end=' ')
\end{verbatim}


\vspace{5mm}
               
{\begin{center}Завдання {5}\end{center}}
\vspace{-6mm}

\begin{verbatim}
def f(n):
    print(n % 10, end=' ')
    if n > 9:
        f(n // 10)
    

f(543216)
\end{verbatim}


\end{minipage}

\newpage

{{22.11.2024 Контрольна робота \No 2, варіант  \No \textbf { 100 }\par}}
{
%\begin {large}

Тестова частина. Завдання 1-5: Що буде виведено за виконання (повідомлення інтерпретатора про помилку позначати error)?

Задачі (на звороті): написати функцію для розв'язання задачі.

\vspace{5mm} \begin{minipage}[t]{8.5cm}
               
{\begin{center}Завдання {1}\end{center}}
\vspace{-6mm}

\begin{verbatim}
a,b,c=5,9,8
def h(b):
    a=2
    b-=1
    c=3
    return a+b+c

a,b,c=3,4,0
print(h(b),a,b,c)
\end{verbatim}


\vspace{5mm}
               
{\begin{center}Завдання {2}\end{center}}
\vspace{-6mm}

\begin{verbatim}
def f():
    try:
        res = int("2")
        return 42
    except Exception: return 4
    except KeyboardInterrupt: return 1
    else: return 34
    finally: return 22
    return res

print(f())
\end{verbatim}


\vspace{5mm}
               
{\begin{center}Завдання {3}\end{center}}
\vspace{-6mm}

\begin{verbatim}
m=5
while m>4:
    m+=-3
print(m)
\end{verbatim}


\end{minipage}
\vline \hspace{6mm}
\begin{minipage}[t]{10.7cm}
               
{\begin{center}Завдання {4}\end{center}}
\vspace{-6mm}

\begin{verbatim}
for c in range(-10,-10+2,-2):
    if c>=4:
        continue
    print(c, end=' ')
    c=7
else:
    print(c, end=' ')
print(c, end=' ')
\end{verbatim}


\vspace{5mm}
               
{\begin{center}Завдання {5}\end{center}}
\vspace{-6mm}

\begin{verbatim}
def f(n):
    print(n % 2, end=' ')
    if n > 2:
        f(n // 2)
    

f(16)
\end{verbatim}


\end{minipage}

\newpage

%end
\end{document}

